% main.tex -- GENERATED by scripts/build_tex.py from main_src.tex; quotations inserted verbatim from ../evidence/corpus.json
\documentclass[11pt,a4paper]{article}
\usepackage{fontspec}
\setmainfont{TeX Gyre Pagella}[Ligatures=TeX]
\setsansfont{TeX Gyre Heros}[Scale=MatchLowercase]
\setmonofont{DejaVu Sans Mono}[Scale=0.85]
\directlua{
  luaotfload.add_fallback("nhifallback", {"NotoSerif:mode=harf;", "NotoSansSymbols:mode=harf;", "DejaVuSans:mode=harf;"})
}
\setmainfont{TeX Gyre Pagella}[Ligatures=TeX, RawFeature={fallback=nhifallback}]
\usepackage[english]{babel}
\usepackage[a4paper,margin=2.3cm]{geometry}
\usepackage{microtype}
\usepackage{graphicx}
\usepackage{booktabs,longtable,array,tabularx}
\usepackage{xcolor}
\usepackage{csquotes}
\usepackage{enumitem}
\usepackage{caption}
\captionsetup{font=small,labelfont=bf}
\usepackage{float}
\usepackage{multicol}
\usepackage[style=authoryear-comp,backend=biber,maxcitenames=2,maxbibnames=6,uniquename=false,uniquelist=false,dashed=false,url=false,doi=true,isbn=false]{biblatex}
\addbibresource{references.bib}
\DeclareFieldFormat{note}{\mkbibbrackets{#1}}
\renewcommand*{\bibfont}{\footnotesize}
\setlength{\bibitemsep}{0.15\baselineskip}
\usepackage[hidelinks]{hyperref}
\hypersetup{pdftitle={Visitors from Above? A Comparative Textual Study of Non-Human Intelligence Motifs in the World's Sacred and Folk Literature},pdfauthor={Investigator in Paranormal Affairs (for Douglas Oliveira)}}
\newcommand{\pid}[1]{\texttt{\small #1}}
\newcommand{\motif}[1]{\textsc{\MakeLowercase{#1}}}
\setlength{\emergencystretch}{2em}
\raggedbottom

\title{Visitors from Above?\\[0.3em]\large A Comparative Textual Study of Non-Human Intelligence Motifs\\ in the World's Sacred and Folk Literature}
\author{Investigator in Paranormal Affairs\\ \small(for Douglas Oliveira)}
\date{26 September 2026}

\begin{document}
\maketitle

\begin{abstract}
\noindent
Do the world's sacred and folk literatures preserve memories of contact with non-human intelligences (NHI) from the sky, the stars or other worlds? This study approaches the question textually and comparatively. Starting from the African section of the Internet Sacred Text Archive (sacred-texts.com), we indexed the whole site (77 categories, 261 entries), filled documented gaps from the aggregator Biblioteca Pleyades (166 items, cited through the original publications it reproduces), and assembled a corpus of 314 script-verified passages from six workstreams: the ancient Near East and Abrahamic traditions, Africa, South and East Asia, Celtic and Northern Europe, the Americas and the Pacific, and modern esoteric and UFO literature as reception. Passages were coded for twelve motifs, typed by source (primary, ethnographic, secondary, fringe, critical) and rated for strength, and five hypotheses were weighed: literal contact, cultural diffusion, shared cognitive myth-making, translation and editorial artefacts, and modern reinterpretation. Sky descent, culture-bringers, divine–human unions, otherworld journeys and hidden peoples are widespread and genuinely old. Forbidden knowledge from heaven is essentially a single Near Eastern literary tradition; supernatural time-lapse clusters in Celtic and Indian sources; no traditional text names a planet of origin; and every vehicle in the primary texts is divine, magical, visionary or a human automaton. The technological reading of ancient beings enters the record between 1896 and 1953, through Theosophy, pulp fiction and the contactees, and it reweights rather than merely rereads the traditional motifs. Texts cannot prove or disprove literal contact; but no attested motif requires it. Diffusion, shared narrative structure, mediation and modern reinterpretation account for the evidence, leaving the universality of the otherworld abduction and time-lapse as genuinely open problems.
\end{abstract}

\vspace{0.5em}
\noindent\textbf{Keywords:} comparative religion; folklore; non-human intelligence; ancient astronauts; Watchers; Dogon; vimāna; fairy abduction; reception history; sacred-texts.com

\tableofcontents
\newpage

%=====================================================================
\section{Introduction and research question}
%=====================================================================

Few modern claims about religion circulate as widely, or with as little attention to the texts on which they purport to rest, as the proposition that the world's scriptures and myths preserve memories of contact with non-human intelligences (NHI): beings from the sky, the stars, other planets or other ``dimensions'' who descended to earth, taught humanity, interbred with it, carried people away, or fought among themselves. The proposition comes in several strengths. In its strongest, ``ancient astronaut'' form, it holds that such passages are more or less literal reports of extraterrestrial visitors misunderstood by pre-scientific observers \parencite{vondaniken1969chariots,sitchin1976planet,temple1976sirius}. In a softer, ``interdimensional'' form associated above all with Jacques Vallée, it holds that fairy lore, angelology and modern abduction reports record one recurrent phenomenon whose nature remains unknown \parencite{vallee1969,vallee1990}. Scholarly responses have mostly been piecemeal: an Assyriologist's rebuttal here, an ethnographic restudy there. Systematic comparison across traditions, organised around the motifs themselves and anchored in verifiable passages, is rarer.

This paper attempts such a comparison. Its question is deliberately modest and textual:

\begin{quote}
\emph{What do the world's religious and mythological texts, as available in standard public-domain translations and ethnographic records, actually say about non-human beings who come from the sky or from other worlds and interact with humankind; how are these motifs distributed across traditions; and which explanatory hypotheses does the evidence, as attested, best support?}
\end{quote}

Five hypotheses were held open throughout, none presupposed:
\begin{description}[leftmargin=2.2em,style=nextline]
\item[H1 Literal contact.] The passages record, however distortedly, real encounters with non-human intelligences.
\item[H2 Cultural diffusion and literary borrowing.] Similarities arise because motifs, stories and texts travelled between cultures.
\item[H3 Shared cognitive and structural myth-making.] Similarities arise independently because human minds and societies generate similar narrative structures (layered cosmoi, sky-fathers, otherworld journeys, hidden peoples).
\item[H4 Translation and editorial artefacts.] Apparent NHI content is produced or amplified by translators, editors, collectors and informant--ethnographer interaction.
\item[H5 Modern reinterpretation.] The NHI reading is supplied by modern esoteric, contactee and ancient-astronaut writers and projected back onto older material.
\end{description}

\noindent One methodological point must be stated at the outset, because it constrains every conclusion that follows. \emph{Texts cannot prove or disprove H1 empirically.} A myth that describes a being descending from heaven is compatible with an actual descent, with a vision, with a poetic image, and with a later interpolation. What textual analysis \emph{can} establish is narrower but still useful: whether the motifs, as they are actually attested, \emph{require} the literal-contact hypothesis in order to be explained---for example, because they contain information unavailable to their authors, or descriptions that resist every other reading---or whether they are adequately accounted for by H2--H5. The paper therefore reports, motif by motif, what the passages say, how old and how mediated they are, and which hypotheses they fit; it does not adjudicate the reality of non-human intelligences in general.

The study began where its commissioner began: at the African section of the online library \emph{Internet Sacred Text Archive} (sacred-texts.com), and expanded from there to the whole site and to the aggregator \emph{Biblioteca Pleyades} (bibliotecapleyades.net). Section~\ref{sec:corpus} describes these sources and the indexing that preceded the close reading; Section~\ref{sec:method} the coding and verification procedure; Section~\ref{sec:results} the regional results, beginning with Africa; Section~\ref{sec:crosscultural} the cross-cultural motif analysis; Section~\ref{sec:reception} the reception history through which the ``ancient NHI'' reading was assembled; and Section~\ref{sec:discussion} the evaluation of H1--H5, the residue of genuinely interesting problems, and the study's limitations.

%=====================================================================
\section{Corpus and sources}\label{sec:corpus}
%=====================================================================

\subsection{The Internet Sacred Text Archive}

The Internet Sacred Text Archive, founded and edited by J.~B. Hare, is a curated library of public-domain translations, ethnographies and esoteric books, organised in top-level categories (African, Native American, Hinduism, Celtic, Theosophy, UFOs, and so on). It is not a scholarly edition: most of its texts are English translations and collections made between roughly 1830 and 1930, many superseded, and the site has a pronounced Anglo-Celtic and Theosophical skew. Its value for this study is twofold. It is where a large public actually meets these texts, and it hosts, side by side, both the primary material and the modern esoteric literature that reinterprets it---which makes it an unusually good laboratory for separating the two.

The indexing phase reviewed all 77 top-level categories of the site (through the archive.sacred-texts.com mirror of the classic index pages) and recorded 261 entries: 253 candidate texts rated for relevance (57 high, 116 medium, 80 low) plus 8 African books rated ``none'' and listed only for completeness, since the African shelf was covered book by book. Of the 77 categories, 35 were judged relevant, 16 cross-listed or marginal, 4 marginal, and 22 without clear relevant material (among them Sikhism, the Bahá'í writings, the Confucian canon, Freemasonry, alchemy, the I Ching and Tarot). Index entries were typed as primary text (P, 87), ethnography or folklore collection (E, 71), secondary scholarship (S, 37), modern esoteric or pseudo-historical (F, 43), modern UFO/contactee/new-religious-movement document (U, 15) and literary fiction (L, 8). Relevance at this stage was judged from tables of contents and descriptions, not full reading; several index descriptions were later corrected in the close-reading phase (Section~\ref{sec:corrections}).

\paragraph{The African entry point.} The African section (\url{https://sacred-texts.com/afr/index.htm}) lists 29 books. It is not a corpus of African scriptures. It combines one medieval literary text, the Ge`ez \emph{Kebra Nagast} in Budge's translation \parencite{budge1922kebra}; ethnographic and folklore collections made between 1868 and 1940 by missionaries, colonial officers and their associates \parencite[e.g.][]{callaway1870religious,theal1886kaffir,dennett1898fjort,ellis1894yoruba,bleeklloyd1911,wyndham1921ife,werner1933bantu,wpa1940drums}; and diasporic and proto-Rastafarian writing. None of the Dogon ethnography that dominates popular discussion of ``African star knowledge'' is hosted there; at indexing, the only Dogon item on the whole site was a 1995 Usenet repost in the UFO files \parencite{dolan1995dogon}.

\subsection{Biblioteca Pleyades as gap-filler}

Biblioteca Pleyades is a large Spanish/English aggregator devoted overwhelmingly to fringe topics. It was used for one purpose: to reach primary and scholarly texts absent from sacred-texts, and the principal fringe works whose claims the study had to test. Its supplementary index lists 166 items (36 high, 81 medium, 49 low relevance; 154 opened, 12 only listed). By type the items are overwhelmingly fringe (F 101), with secondary commentary (S 32), primary texts (P 23), sceptical critique (C 6) and ethnography (E 4). Nine gaps in the sacred-texts coverage were defined and checked (Table~\ref{tab:gaps}). Wherever an aggregator page reproduced a published work, the study cites the original publication and treats the aggregator only as an access copy; aggregator pages that reproduce copyrighted material, carry unattributed translations, or insert editorial remarks into supposedly original text are flagged as such.

\begin{table}[ht]
\centering\small
\caption{Gaps in sacred-texts.com coverage and their status after the Pleyades gap-fill (from \texttt{index\_pleyades.md}).}\label{tab:gaps}
\begin{tabularx}{\textwidth}{@{}p{3.9cm}p{2.1cm}X@{}}
\toprule
Gap & Status & What was gained / what is still missing \\
\midrule
G1 Dogon, Nommo, Sirius & partial & Temple's book, including an English translation of \textcite{griaule1950sirius}; the critique only via Coppens's summary of van Beek. \\
G2 Vaimanika Shastra, Samarāṅgaṇa & partial & Josyer's 1973 translation and the 2001 AR\&DB study; the 1974 IISc critique fetched directly from its authors' site. \\
G3 Sitchin and critiques & filled & Sitchin's books, Heiser's philological critique, and Dalley's Atrahasis and Adapa. \\
G4 East/Southern Africa & partial & Credo Mutwa's 1999 interview and essays (fringe-mediated); no Zulu or Swahili ethnography. \\
G5 Mesoamerica & partial (primary filled) & Full Goetz--Morley \emph{Popol Vuh}; Quetzalcoatl material fringe only. \\
G6 Andes & partial & Sarmiento duplicates sacred-texts; no Huarochirí manuscript. \\
G7 Australia/Pacific & partial (weak) & One inaccurate Spanish Wandjina page; Solomon Islands UFO lore. \\
G8 China/Tibet & partial & No Chinese primary texts; Korean foundation myths (mislabelled ``China''); Berzin on Shambhala. \\
G9 Primary texts not on sacred-texts & partial-to-filled & Qumran Book of Giants, Genesis Apocryphon, Lockett's Hopi ethnography, Korean texts. \\
\bottomrule
\end{tabularx}
\end{table}

\subsection{Six regional workstreams and the verified corpus}

Close reading was organised in six workstreams: the ancient Near East and Abrahamic traditions (\texttt{neareast}); Africa (\texttt{africa}); South and East Asia (\texttt{asia}); Celtic and Northern Europe (\texttt{celtic}); the Americas and the Pacific (\texttt{americas\_pacific}); and modern esoteric and UFO literature as \emph{reception} (\texttt{modern}). Each produced a file of verified passages, an analytic synthesis, and a bibliography. Merged, the corpus contains \textbf{314 verified passages} (Table~\ref{tab:corpus}); a further 15 verified African passages cut for length are preserved separately.

\begin{table}[ht]
\centering\small
\caption{The verified passage corpus by area, source type and strength rating (\texttt{evidence/corpus.json}).}\label{tab:corpus}
\begin{tabular}{@{}lrrrrrrrrr@{}}
\toprule
Area & $n$ & P & E & S & F & C & strong & moderate & weak \\
\midrule
Near East \& Abrahamic & 57 & 45 & 0 & 9 & 2 & 1 & 35 & 19 & 3 \\
Africa & 50 & 7 & 19 & 11 & 8 & 5 & 17 & 13 & 20 \\
South \& East Asia & 51 & 37 & 1 & 5 & 2 & 6 & 29 & 19 & 3 \\
Celtic \& N. Europe & 52 & 26 & 15 & 5 & 6 & 0 & 24 & 22 & 6 \\
Americas \& Pacific & 50 & 8 & 29 & 4 & 7 & 2 & 27 & 13 & 10 \\
Modern reception & 54 & 3 & 0 & 3 & 47 & 1 & 41 & 8 & 5 \\
\midrule
Total & 314 & 126 & 64 & 37 & 72 & 15 & 173 & 94 & 47 \\
\bottomrule
\end{tabular}
\end{table}

\subsection{Source-type classification}

Every passage carries one of five source-type codes, which matter more for the argument than any other variable:
\begin{description}[leftmargin=1.5em]
\item[P] \emph{Primary text in translation}: a scripture, epic, chronicle, medieval literary text, ballad, or (in the modern area) a primary document of modern religion or fiction such as Lovecraft or the \emph{Vaimanika Shastra}.
\item[E] \emph{Ethnographic record}: texts collected from named or unnamed informants by outsiders, c.~1690--1980.
\item[S] \emph{Secondary scholarship} and compilations (e.g. Werner, Ginzberg, Kramer's paraphrases, translators' notes).
\item[F] \emph{Fringe, esoteric or new-revelation} literature: Theosophy, channelled scripture, contactee books, ancient-astronaut writing, BBS files.
\item[C] \emph{Critical scholarship} evaluating a specific fringe claim (e.g. van Beek, Heiser, Mukunda et al.).
\end{description}
F passages are primary evidence \emph{for the history of modern religion}, not for the beliefs of the traditions they reinterpret. Figure~\ref{fig:types} shows the composition by area: the Near East and Asia are dominated by primary texts, Africa and the Americas by ethnographic records, and the modern area by fringe literature by design.

\begin{figure}[ht]
\centering
\includegraphics[width=0.85\textwidth]{figures/fig3_source_types.pdf}
\caption{Source-type composition of the 314 verified passages by area.}\label{fig:types}
\end{figure}

\subsection{Coverage and gaps}

Absence from sacred-texts is not absence from a tradition, and several important gaps remained after the gap-fill. Africa is represented mostly by southern and western material and the diaspora; Ellis's Tshi and Ewe volumes are not hosted; there is no Dogon ethnography on the site and no Zulu ethnography on the aggregator. China is represented only by a modern hoax (the ``Dropa stones'') and by Werner's survey, which was indexed but not excerpted; no Chinese primary text (Shanhaijing, Huainanzi, Chuci) was read. \emph{Lebor Gabála Érenn} is not on sacred-texts in full and was read in a Macalister excerpt from another site. The sacred-texts \emph{Nihongi} is excerpts only. For Mesoamerica, no Sahagún, Anales de Cuauhtitlan or Chilam Balam was fetched, so the pre-Hispanic basis of the ``returning Quetzalcoatl'' could not be tested directly. These gaps bound the conclusions: they are conclusions about the attested, English-accessible corpus.

%=====================================================================
\section{Method}\label{sec:method}
%=====================================================================

\subsection{Motif coding}

Passages were coded for twelve motifs defined during indexing (Appendix~\ref{app:codebook}): sky descent (\textsc{sky}), culture-bringer (\textsc{cult}), forbidden-knowledge transfer (\textsc{forb}), hybrid offspring and divine--human unions (\textsc{hyb}), celestial vehicles (\textsc{veh}), star-origin ancestry and star-people (\textsc{star}), otherworld time dilation (\textsc{time}), abduction or otherworld journey (\textsc{abd}), plurality of inhabited worlds (\textsc{plur}), other planes and hidden peoples (\textsc{dim}), channelled or contactee communication (\textsc{chan}), and war or rebellion among sky beings (\textsc{war}). A passage may carry several codes. The codes are \emph{etic} and descriptive: they record that a passage describes, say, a being descending from the sky, not that the being is an extraterrestrial. The emic identity---angel, Watcher, orisha, deva, kami, jinn, \emph{síd}-folk, kachina---was recorded alongside in the passage's notes, following the recommendation of the Near Eastern workstream that motif coding must not collapse emic categories.

\subsection{Passage verification}

Each quotation in the corpus was verified by script against the text of the page from which it was taken, after whitespace normalisation, on 26 September 2026. Quotations were kept short (the Asian workstream set a 60-word ceiling), with ellipses marking omissions. Each record gives the tradition, text title, translator and date, URL, a locator (chapter, section or page), the verbatim quotation, the motif codes, a summary, interpretive notes and a strength rating. All quotations reproduced in this paper were inserted mechanically from the merged corpus file, so that no quotation in the paper differs from its verified source string.

\subsection{Strength rating}

``Strength'' rates how clearly a passage exhibits the motif(s) coded, not how credible or important it is. A strong \textsc{veh} passage describes an aerial vehicle unambiguously; a weak one requires interpretation. In the modern area the rating describes how clearly the motif appears \emph{in the modern text}; it is never a judgement that the ancient tradition held the idea. Overall, 173 passages are rated strong, 94 moderate and 47 weak (Figure~\ref{fig:strength} in Appendix~\ref{app:data}).

\subsection{Dating, provenance and mediation}

Because the diffusion and contact hypotheses make different predictions about dates and routes, each workstream recorded the date of composition, the date and nature of the edition used, and the chain of mediation (collector, translator, editor, aggregator). Where a translation is superseded, the notes name the modern standard edition (e.g. ETCSL for Sumerian literature, \textcite{lambert2013creation} for Enūma eliš, \textcite{izreel2001adapa} for Adapa, \textcite{tedlock1996popol} for the \emph{Popol Vuh}), even when the older public-domain version was the one read.

\subsection{Evaluating the hypotheses}

For each motif and region the analysis asked what each hypothesis predicts. H1 predicts content that exceeds its authors' knowledge (astronomical, technical or biological information not otherwise available), physical descriptions that resist every non-literal reading, and, ideally, independent attestations of the same visitors. H2 predicts shared names, plots and specific details along plausible routes of contact, with dates consistent with transmission. H3 predicts structural similarity without shared specific detail, broad geographical distribution, and association with universal experiences and social functions. H4 predicts that the NHI-like feature appears in a particular translation or collection but not in the source language or in other recordings. H5 predicts that the NHI reading has a datable first appearance in modern literature and a documented chain of borrowing.

\subsection{Corrections to the index}\label{sec:corrections}

Close reading corrected several index descriptions; these corrections are applied throughout the paper and in the revised \texttt{index.md} (the original is kept as a backup):
\begin{itemize}[itemsep=1pt]
\item The sacred-texts \emph{Cath Maige Tuired} page (\texttt{neu/cmt/cmteng.htm}) reproduces Elizabeth A. Gray's 1982 Irish Texts Society translation \parencite{gray1982cmt}, not Whitley Stokes (1891).
\item Alice Bailey's \emph{Initiation, Human and Solar}, chapter IV, dates Sanat Kumara's descent to ``approximately eighteen million years ago'' but never names Venus \parencite{bailey1922initiation}; the Venus origin belongs to Scott-Elliot, Leadbeater and Ballard.
\item The chapters of Weil's \emph{Biblical Legends of the Mussulmans} fetched from sacred-texts (Adam, Idris, Solomon) do not contain the Hārūt and Mārūt narrative; that tale is Qur'ān 2:102 plus exegetical commentary.
\item Ogumefu's ``Kingdom of the Yorubas'' contains no sky descent \parencite{ogumefu1929yoruba}.
\item The \emph{Kebra Nagast}'s ``chariot'' chapters (113--117) describe no flight.
\item \emph{Chariots of the Gods?} appeared in German in 1968 and in English in 1969 \parencite{daniken1968erinnerungen,vondaniken1969chariots}, not 1971.
\item The sacred-texts \emph{Popol Vuh} pages now return 404; the Goetz--Morley English translation of Recinos, reproduced on Pleyades, was used \parencite{goetz1950popol}.
\item Vallée's quotation of Evans-Wentz, as reproduced online, differs from the 1911 text and must be checked against the print edition (Section~\ref{sec:magonia}).
\end{itemize}

%=====================================================================
\section{Results by region}\label{sec:results}
%=====================================================================

The regional sections below follow a common pattern: sources and their mediation, the motifs genuinely present, the key passages (quoted verbatim, identified by corpus ID), what belongs to modern reinterpretation, and a summary. Africa comes first, because it was the entry point of the study and because its material illustrates with particular clarity every one of the five hypotheses.

%---------------------------------------------------------------------
\subsection{Africa}\label{sec:africa}
%---------------------------------------------------------------------

\subsubsection{Sources and mediation}

The African workstream read all 29 books on the sacred-texts African shelf, plus the Dogon and Credo Mutwa material reproduced on Pleyades and the main scholarly critiques; it retained 50 passages (Table~\ref{tab:corpus}). Mediation is the governing caveat. Almost every oral tradition here was collected under colonial conditions and translated by Europeans. Callaway, an Anglican missionary, printed transcripts that expose interviewer effects: in one, an old woman's account of Unkulunkulu moves from ``beneath'' to ``above'' after leading questions, and her son objects: \enquote{But Ubapa said, "No! she now begins to speak at cross purposes. She did not say this to the Missionary yesterday. She said Unkulunkulu was from beneath. But now she says he was from above."}~\pid{[AFR-RSA-03]} \parencite{callaway1870religious}. Wyndham, an Assistant District Officer, turned what he called ``bare and uncertain'' Ife legends into blank verse on Ford Madox Hueffer's advice and after reading Frazer \parencite{wyndham1921ife}. Werner's \emph{Myths and Legends of the Bantu} is a secondary compilation drawing mostly on German Lutheran, Swiss and Catholic missionaries, still framed by the discredited ``Hamitic'' hypothesis \parencite{werner1933bantu}. Bleek and Lloyd's /Xam narrators included men held as convicts at the Breakwater prison in Cape Town \parencite{bleeklloyd1911}. None of this makes the material worthless, but it obliges us to speak of \emph{documented colonial-era versions}, not of ``ancient African teachings''.

\subsubsection{The populated sky}

Sky descent is the strongest and most widespread African motif. The Chaga (Wachaga) material in Werner is the clearest case of an ancestral sky-race: 
\begin{quote}\small One clan of the Wachaga claims that its ancestor fell from the sky during a rainstorm. He belonged to a race called the Wakyambi, living in the sky, "far above the sun," and having tails.\par\hfill{\footnotesize --- Myths and Legends of the Bantu, ch. V, "Tailed Heaven-folk" (Chaga) \pid{[AFR-MLB-08]}}\end{quote}
 A tailed couple descends and explains, \enquote{When asked where they came from they answered, "God has sent us down on a cloud. We are looking for a place to live in."}~\pid{[AFR-MLB-09]} In Yoruba tradition the gods descend from Heaven by a chain to create land at Ife: \enquote{On reaching the edge of Heaven, Odúwa hung a chain over the cliff and sent down a priest, called Ojúmu, with a snail-shell full of magic sand and a "five-fingered" fowl.}~\pid{[AFR-IFE-01]} \parencite{wyndham1921ife}; in Ellis's version Obatala saves survivors of a flood by drawing them \enquote{up into the sky by means of a long iron chain}~\pid{[AFR-YOR-01]} \parencite{ellis1894yoruba}. Callaway's elders located \enquote{a nation of people}~\pid{[AFR-RSA-02]} in heaven with the Creator.

The single most ``visitor-like'' passage in the whole African corpus is also from Callaway:

\begin{quote}\small one Unkulunkulu came from beneath; and another descended from above in a fog. They did not understand him who came down in a fog. They say he was altogether white.\par\hfill{\footnotesize --- The Religious System of the Amazulu, Unkulunkulu, p. 39 (Ushunguwane Zimase) \pid{[AFR-RSA-01]}}\end{quote}

In the fuller text the white being says he is a man, refuses local food, and vanishes in another fog. The African workstream identified at least three readings. It may be a sky-ancestor epiphany in mist. It may be a mythologised memory of early white contact, a reading supported by the fact that Callaway's informants repeatedly equate whites with sky-lords, and that Ife tradition made a son of Orisha the ancestor of the white races (a further verified passage kept in the full African file). Or, least economically, it may record a literal visitation. The text itself cannot decide among these; but the second reading has independent ethnographic support that the third lacks.

What the African sky is \emph{like} matters as much as who descends from it. Across the corpus the sky is a solid roof over a mirror society. Werner summarises the means of access: \enquote{Either they climb a tree, or they ascend by means of a rope,[1] or the spider obligingly spins a thread for them.}~\pid{[AFR-MLB-06]} A heroine \enquote{soon found herself in the country above the sky, which appeared to be not unlike the one she had left.}~\pid{[AFR-MLB-07]} In Dennett's Kongo material the spider \enquote{reached the blue roof of heaven}~\pid{[AFR-FJO-01]} and left a thread by which others climb \parencite{dennett1898fjort}; a Chaga boy commands a magic stool, \enquote{Then he said, "Stool, go up on high, as my father's rope does when he hangs up his beehive in the forest!"}~\pid{[AFR-MLB-11]} The heaven-dwellers keep villages, gardens and cattle. Many sky beings are simply weather personified: the Thunder appears \enquote{in the likeness of a man, with a little bright axe in his hand.}~\pid{[AFR-MLB-03]} This is a cosmology of adjacent, ladder-reachable worlds, not of distant space.

\subsubsection{Culture-bringers, and a warning from a footnote}

The descended Yoruba gods \enquote{began to teach their arts and crafts to men.}~\pid{[AFR-IFE-02]} The Chaga Murile story reverses the pattern: a human visitor to the Moon-chief's country \enquote{found the people there eating their food raw}~\pid{[AFR-MLB-12]}, because they did not know fire. Mary Kingsley's editorial footnote to Dennett is methodologically decisive for the whole culture-bringer category. A Kongo river-spirit credited with teaching the planting of bananas and manioc is being credited with crops that, she observes, \enquote{Bananas and manioc were introduced in Fjort culture by the Roman Catholic missionaries, who first landed among the Fjort in 1490}~\pid{[AFR-FJO-02]} \parencite{dennett1898fjort}. Culture-bringer attributions update themselves to absorb recent imports; they cannot be read as records of where the arts actually came from.

\subsubsection{Sky marriage, abduction and hybrid offspring}

The Rwandan ``Thunder's Bride'' is a full sky-marriage: \enquote{The Thunder caught her up and carried her off to the sky and married her.}~\pid{[AFR-MLB-04]} She bears three children, and later she and her party are caught up into the air. Kintu's marriage to the daughter of Gulu (Heaven) among the Ganda, the Ila tower to heaven, and the Chaga visit to the sun-king's shining court belong to the same field. The Gullah ``flying Africans'' of the WPA interviews, \enquote{Lots uh slabes wut wuz brung obuh frum Africa could fly.}~\pid{[AFR-DAS-01]} \parencite{wpa1940drums}, are a resistance legend, not contact lore. All of these belong to pan-African and worldwide tale-types (sky-marriage, tower to heaven, false bride), which favours shared narrative structure (H3) over any single historical event.

\subsubsection{The Kebra Nagast: flying wagons and Ethiopic Watchers}

The \emph{Kebra Nagast}, compiled in Ge`ez in the early fourteenth century to legitimise the Solomonic dynasty restored in 1270, supplies Africa's most cited ``vehicle''. Among Solomon's gifts to the Queen of Sheba is \enquote{a vessel wherein one could travel over the sea, and a vessel wherein one could traverse the air (or winds), which Solomon had made by the wisdom that God had given unto him.}~\pid{[AFR-KN-01]} When Menyelek carries the Ark to Ethiopia, the Archangel Michael leads the caravan: 
\begin{quote}\small no man hauled his wagon, but he himself (i.e. Michael) marched with the wagons, and whether it was men, or horses, or mules, or loaded camels, each was raised above the ground to the height of a cubit\par\hfill{\footnotesize --- Kebra Nagast, ch. 52, Budge pp. 77-78 \pid{[AFR-KN-02]}}\end{quote}
 Witnesses report that the travellers went \enquote{in wagons that were suspended in the air; and they were swifter than the eagles that are in the sky}~\pid{[AFR-KN-03]} \parencite{budge1922kebra}. In context these are miracles of divine favour legitimising the translation of the Ark, and they sit squarely in the Solomon-and-the-wind legend cycle known from the Qur'ān (21:81, 34:12, 38:36) and Jewish flying-carpet legends. The ``chariot'' of chapters 113--117, sometimes cited as a further vehicle, is a royal palladium: \enquote{And the chariot belongeth to the King of Ethiopia, and it hath vanquished his enemy.}~\pid{[AFR-KN-09]}

Chapter 100 is an Ethiopic variant of the Watchers myth discussed in Section~\ref{sec:neareast}. Angels who mocked Adam are given bodies: 
\begin{quote}\small And straightway there were given unto them with His word flesh, and blood, and a heart of the children of men. And they were content to leave the height of heaven, and they came down to earth\par\hfill{\footnotesize --- Kebra Nagast, ch. 100, Budge p. 187 \pid{[AFR-KN-05]}}\end{quote}
 The daughters of Cain whom they approach conceive giants who cannot be born naturally: \enquote{the daughters of Cain with whom the angels had companied conceived, but they were unable to bring forth their children, and they died.}~\pid{[AFR-KN-06]} Unlike 1~Enoch, however, the angels here teach nothing; artisanship is attributed to the Cainites. Forbidden-knowledge transfer from sky-beings (\textsc{forb}) is therefore essentially absent from the African corpus. The Ife ``War of the Gods'' over a stolen bag of arts, which \enquote{is said to have lasted 201 years}~\pid{[AFR-IFE-03]}, is the only indigenous war among sky-beings and reads naturally as a dynastic charter.

\subsubsection{Stars: a /Xam control case}

Authentic African star lore on sacred-texts is /Xam. A girl of the Early Race makes the Milky Way from wood ashes: \enquote{she put her hands into the wood ashes; she threw up the wood ashes into the sky. She said to the wood ashes: "The wood ashes which are here, they must altogether become the Milky Way.}~\pid{[AFR-SBF-01]} Stars are persons; Sirius and Canopus are female ``grandmother'' stars, addressed with burning sticks: \enquote{He points (it) burning towards Sirius; he says that Sirius shall twinkle like Canopus.}~\pid{[AFR-SBF-03]} \parencite{bleeklloyd1911}. This is an important control. It is detailed Sirius lore, recorded in the 1870s from a southern African people, with no invisible companion star and no visitors.

\subsubsection{The Dogon ``Sirius mystery''}

Griaule and Dieterlen's 1950 article reported that initiates regard Sirius as \enquote{one of the foci of the orbit of a tiny star called Digitaria, po tolo}~\pid{[AFR-DOG-01]}, the smallest and heaviest star: \enquote{Because of this role, the star which is considered to be the smallest thing in the sky is also the heaviest: 'Digitaria is the smallest thing there is. It is the heaviest star:'}~\pid{[AFR-DOG-02]} \parencite{griaule1950sirius}. The article makes no claim of alien contact; Nommo appears in it as \enquote{Instructor of the new world}~\pid{[AFR-DOG-05]}. It was \textcite{temple1976sirius} who turned Nommo into \enquote{the great culture-hero and founder of civilization who came from the Sirius system to set up society on the Earth.}~\pid{[AFR-DOG-06]} According to Coppens, summarising Picknett and Prince, Temple received the article from Arthur M. Young, a participant in the channelled ``Council of Nine'' circle \parencite{coppens2000sirius,picknettprince1999stargate}.

Three lines of critique bear on the claim. First, the 1950 data do not fit Sirius~B well. The article itself says: \enquote{When Digitaria is close to Sirius, the latter becomes brighter; when it is at its most distant from Sirius, Digitaria gives off a twinkling effect, suggesting several stars to the observer.}~\pid{[AFR-DOG-03]} That describes a visible object, which Sirius~B is not \parencite{jamesthorpe1999mysteries}; and ``smallest/heaviest'' sits inside a seed cosmogony in which \emph{po} is fonio, the smallest grain. Second, contamination before Griaule is possible: \textcite{ridpath1978sirius}, following \textcite{pesch1977dogon} and \textcite{sagan1979broca}, noted that \enquote{Peter and Roland Pesch of the Warner and Swasey Observatory in Ohio have pointed out that French schools have existed in the Dogon area since 1907.}~\pid{[AFR-DOG-12]} Ridpath concedes that no document shows such transmission occurred. Third, and most important, van Beek's decade-long restudy found that the Dogon know Sirius as \enquote{dana tolo, the hunter's star}~\pid{[AFR-DOG-09]} and that \enquote{Knowledge of the stars is not important either in daily life or in ritual.}~\pid{[AFR-DOG-09]} No Dogon outside Griaule's informant circle knew \emph{sigu tolo} or \emph{po tolo}; van Beek concluded that it was \enquote{Griaule himself, driven by his own convictions, who transformed po tolo and sigu tolo into a mystery}~\pid{[AFR-DOG-10]} \parencite{vanbeek1991dogon}. Calame-Griaule replied that van Beek used different informants and methods \parencite{calamegriaule1991}, and \textcite{apter2005griaule} argued that esoteric ``deep knowledge'' in West Africa is fluid and context-bound, so its non-replication does not prove fabrication. Even on Apter's reading, the Sirius system is the product of a particular informant--ethnographer exchange in the 1940s, not a fixed ancient doctrine. The Dogon case therefore belongs to the history of anthropology (H4) and of modern reinterpretation (H5), not to the evidence for contact.

\subsubsection{Credo Mutwa and the Chitauri}

Vusamazulu Credo Mutwa (c.~1921--2020) is a modern informant, not a transmitter of documented Zulu tradition. In a 1999 interview arranged through David Icke he glossed \emph{izulu} as \enquote{inter-planetary space}~\pid{[AFR-MUT-01]}, contrary to the standard glosses sky, heaven, weather; described the Chitauli as \enquote{the dictators, the ones who tell us the law}~\pid{[AFR-MUT-02]}; and identified the Mantindane with modern abduction lore: \enquote{We believe, sir, that the Mantindane ("the tormentor"), the Greys , are really servants of the Chitauli .}~\pid{[AFR-MUT-03]} \parencite{martin1999mutwa}. His retelling of Nommo arriving \enquote{in their fantastic sky ship}~\pid{[AFR-MUT-07]} follows Temple \parencite{mutwa_web}. The clearest diagnostic is the \emph{imbulu}. Theal (1886) glossed it as follows: \enquote{The mbulu is a fabulous creature, firmly believed in by little folks. It can assume the human form, but cannot part with its tail.}~\pid{[AFR-XFT-01]} \parencite{theal1886kaffir}. Werner places it in the Zulu false-bride tale. Mutwa turns the same figure into a reptilian non-human: \enquote{she met a creature called an Imbulu. And this Imbulu was a lizard which has the body and the limbs of a human being, but a long tail.}~\pid{[AFR-MUT-04]} ``Chitauri'' appears nowhere in the nineteenth-century Zulu material on sacred-texts. \textcite{chidester2002mutwa} treats Mutwa's authenticity as invented and appropriated; \textcite{podolecka2023mutwa}, whose 26 sangoma informants did not recognise his mythic figures, concluded: \enquote{Hence Mutwa’s stories – even though myths in their construction, narrative and aim – are not Zulu myths, and they do not convey Zulu beliefs about their origins.}~\pid{[AFR-POD-01]} Mutwa is best used as primary evidence for contemporary contactee religion and its globalisation---a feedback loop in which a Western fringe reading returns as ``African tradition''.

\subsubsection{Negative findings and summary}

A table-of-contents review plus sampled chapters found no NHI-relevant motifs in Theal (apart from the mbulu), Talbot, Hearn, Du Bois, James's \emph{Stolen Legacy}, Houston, \emph{The Promised Key}, \emph{The Wisdom of Rastafari} and Macdonald; only marginal material in Honey, Dayrell, Rattray, Beckwith, Williams, Nassau and Dennett's \emph{Black Man's Mind}. Dennett's ``Ngomba's Balloon'' involves no flight in the portion read; Ogumefu's Yoruba kingdom has no sky descent: \enquote{The resourceful Oranyan spread upon the water his pieces of iron, and upon the iron he placed the scrap of cloth, and upon the cloth the soil, and on the soil the cock.}~\pid{[AFR-YL-01]} Across all 29 books there is no time dilation, no channelled contact before Mutwa, and no forbidden-knowledge transfer from sky-beings.

\emph{Summary.} Genuinely present: sky-descending gods and ancestors (strong in Yoruba and Chaga material), a populated sky reached by ladder-like means, sky marriages with offspring, miraculous aerial conveyance within Ethiopian Christian-Solomonic legend, and an Ethiopic Watchers tradition derived from Jewish apocalyptic. Modern reinterpretation: the extraterrestrial Nommo, the reptilian Chitauri and hybridising Greys, and ``flying machines'' in the \emph{Kebra Nagast}. No passage requires literal contact.

%---------------------------------------------------------------------
\subsection{The ancient Near East and the Abrahamic traditions}\label{sec:neareast}
%---------------------------------------------------------------------

\subsubsection{Sources}

This area holds the richest single cluster of sky-being material in the study (57 passages, 45 of them primary). It has a Mesopotamian layer---Sumerian and Akkadian myths of the Anunnaki and Igigi, the \emph{me}, Adapa, Etana and Atrahasis, and Berossus on Oannes---and a Jewish, Christian, Manichaean and Islamic layer organised around the reception of Genesis 6:1--4: the Enochic corpus, Jubilees, the Aramaic Book of Giants, the Genesis Apocryphon, 2~Enoch, Manichaean fragments, rabbinic and medieval legend (through Ginzberg), and Qur'ānic material on jinn, Hārūt and Mārūt, and Idrīs. The dates range from an Old Babylonian Atrahasis copy of the seventeenth century BCE to the eleventh-century Midrash of Shemhazai and Azael; the Book of Watchers is third century BCE, the Similitudes late first century BCE to first century CE \parencite{charles1917enoch,nickelsburg2001enoch,drawnel2011astronomical}. Several translations read are superseded (Kramer 1944, King 1902, Rogers 1912, Mackenzie 1915), and the aggregator copies of Dalley, of Wise, Abegg and Cook, and of an unattributed Genesis Apocryphon must be checked against print \parencite{dalley1989myths,wise1996dss,vermes2004complete}.

\subsubsection{Sky descent and hybrid offspring}

These dominant motifs are genuinely in the texts. The seed is Genesis 6:

\begin{quote}\small That the sons of God saw the daughters of men that they were fair; and they took them wives of all which they chose. ... There were giants in the earth in those days; and also after that, when the sons of God came in unto the daughters of men\par\hfill{\footnotesize --- Genesis 6, Gen 6:2, 6:4 \pid{[NE26]}}\end{quote}

The Book of Watchers expands it into a descent of two hundred: \enquote{And the angels, the children of the heaven, saw and lusted after them, and said to one another: 'Come, let us choose us wives from among the children of men and beget us children.' ... And they were in all two hundred; who descended}~\pid{[NE01]} \parencite{charles1917enoch}. The motif recurs in Jubilees (NE18), 2~Enoch (NE23), the Qumran and Manichaean Books of Giants (NE30--33), and in the Genesis Apocryphon and 1~Enoch 106, which show a live anxiety about suspected angelic paternity. Lamech says of the newborn Noah: \enquote{I have begotten a strange son, diverse from and unlike man, and resembling the sons of the God of heaven; and his nature is different and he is not like us ... it seems to me that he is not sprung from me but from the angels}~\pid{[NE16]} Charles's note to Jubilees 5:1 records that the angelic reading of Genesis 6 ``represents the older Jewish exegesis, which was later given up'' \parencite{charles1917jubilees}: the divine-being reading was the ancient mainstream, not a modern invention, and the rabbinic ``sons of judges'' and patristic Sethite readings are later apologetic re-readings.

\subsubsection{Culture-bringers and forbidden knowledge: a matched pair}

Mesopotamian sources present civilisation as a positive gift from the Abzu. Berossus's Oannes emerges from the sea and \enquote{taught them to construct cities, to found temples, to compile laws}~\pid{[NE39]} \parencite{cory1832fragments}; Enki holds \enquote{all the divine decrees that are fundamental to civilization}~\pid{[NE41]} \parencite{kramer1944sumerian}. The Book of Watchers inverts this. Its angels transmit crafts illicitly:

\begin{quote}\small And Azâzêl taught men to make swords, and knives, and shields, and breastplates, and made known to them the metals ... Barâqîjâl, ... (taught) astrology, Kôkabêl the constellations, Ezêqêêl the knowledge of the clouds\par\hfill{\footnotesize --- 1 Enoch, 1 En 8:1-3 \pid{[NE03]}}\end{quote}

The archangels charge that Azazel \enquote{revealed the eternal secrets which were (preserved) in heaven, which men were striving to learn}~\pid{[NE04]}. The Similitudes even make writing a fallen gift: Penemue \enquote{instructed mankind in writing with ink and paper}~\pid{[NE13]}. Jubilees softens the inversion, saying the Watchers were originally sent \enquote{that they should instruct the children of men}~\pid{[NE17]}, and has their teaching survive the flood on carved rock (NE19), a Mesopotamian scribal topos. The closest structural parallel to 1~Enoch 9:6 is Anu's complaint in Adapa: \enquote{Why has Ea revealed to impure mankind The heart of heaven and earth?}~\pid{[NE46]} \parencite{rogers1912cuneiform}. The Qur'ān's Hārūt and Mārūt receive magic that \enquote{came down at Babylon}~\pid{[NE49]}, and teach it only as a trial \parencite{yusufali1934quran}---the Babylonian setting is itself telling.

\subsubsection{Ascent and return, vehicles, stars and war}

Enoch is the paradigm of ascent. The laconic seed---\enquote{And Enoch walked with God: and he was not; for God took him.}~\pid{[NE27]}---grows into the dream-ascent to a crystal house and wheeled throne (NE07), the cosmic tours with Uriel, the two shining giants at the bedside in 2~Enoch (NE21), angelic re-clothing (NE24), and a timed stay: \enquote{He was taken up to heaven on the first day of the month Tsivan and remained in heaven sixty days.}~\pid{[NE25]} \parencite{morfill1896secrets}. Sixty days in heaven correspond to sixty days on earth: a negative finding for time dilation. The ancient texts consistently frame these journeys as visions, summonses or divine transfers; none describes capture by an agent hostile to the human subject.

The vehicle motif is weaker than its popular reputation. Ezekiel's vision is explicitly analogical:

\begin{quote}\small a whirlwind came out of the north, a great cloud, and a fire infolding itself, and a brightness was about it ... their appearance and their work was as it were a wheel in the middle of a wheel. ... when the living creatures were lifted up from the earth, the wheels were lifted up\par\hfill{\footnotesize --- Ezekiel 1, Ezek 1:4-5, 1:16, 1:19 \pid{[NE28]}}\end{quote}

The prophet identifies it as ``the likeness of the glory of the LORD'' and, in chapter 10, as cherubim: \enquote{behold the four wheels by the cherubims}~\pid{[NE29]} \parencite{kjv1611}. In 1~Enoch 14 the wheels become God's throne-chariot; the sun's ``chariots'' in the Astronomical Book and 2~Enoch are conventional solar imagery. Platt's 1926 footnote to 2~Enoch 11, \enquote{Cf. "Rapid Transit."}~\pid{[NE22]}, is a neat, datable example of an editor modernising the imagery \parencite{platt1926forgotten}.

In 1~Enoch stars are sentient beings who can transgress and be imprisoned: \enquote{the stars which roll over the fire are they which have transgressed the commandment of the Lord}~\pid{[NE11]}. In the Animal Apocalypse a fallen star is the conventional code for an angel (NE14). This is the ancient star--angel equation, not ``star-people ancestry''. Rebellion among divine beings is robust: the Igigi strike in Atrahasis, \enquote{Every single one of us declared war! We have put a stop to the digging.}~\pid{[NE43]} \parencite{dalley1989myths}, the binding of the Watchers, the Grigori rebellion of 2~Enoch, and Iblis.

\subsubsection{Other orders of being}

1~Enoch 15 gives a formal ontology: the giants' disembodied spirits become demons on earth (NE08). The jinn are a third intelligent order \enquote{From fire free of smoke}~\pid{[NE51]}, able to eavesdrop on heaven: \enquote{we pried into The secrets of heaven; But we found it filled With stern guards And flaming fires.}~\pid{[NE50]} Weil's ``kingdom of spirits'' \enquote{fills up the whole space between the earth and heaven}~\pid{[NE53]} \parencite{weil1846bible}. The only real plurality-of-worlds item is the midrashic seven earths (NE38), nested strata each with ``its own sky'', not planets \parencite{ginzberg1909legends}.

\subsubsection{Sitchin tested against the texts}

Sitchin's \emph{The 12th Planet} (1976) made four core claims: that \emph{shem} and Sumerian MU mean ``rocket'', so that Genesis 6:4's ``men of renown'' are ``people of the rocket ships'' (\enquote{the term shem must be taken in its original meaning - a rocket, a rocket ship}~\pid{[NE54]}); that the Anunnaki were astronauts from a twelfth planet, Nibiru; that the Igigi revolt was a miners' strike; and that humans were genetically engineered as a ``Primitive Worker'' \parencite{sitchin1976planet}. Tested against the primary texts, the first claim fails philologically. Heiser shows that the lexical lists equate MU with \emph{šamû} (``heaven'', ``rain''), that Akkadian \emph{šumu} and Hebrew \emph{šem} mean ``name'', and that Sumerian lacks the relative pronoun Sitchin's etymology requires; his verdict is that \enquote{the words Mr. Sitchin tells us refer to rocket ships have no such meaning according to the ancient Mesopotamians themselves.}~\pid{[NE55]} \parencite{heiser_sitchin_rockets}. The second is unsupported: \emph{nēberu} is a ``crossing'' point applied to Marduk's star, Jupiter or Mercury \parencite{heiser_nibiru,horowitz1998cosmic}. The third and fourth have a real textual kernel with a transformed meaning. Atrahasis does narrate a strike of junior gods against forced digging, and humanity's creation from clay mixed with a slain god's flesh and blood: \enquote{Nintu shall mix the clay With his flesh and blood. Then a god and a man Will be mixed together in clay.}~\pid{[NE44]} Enūma eliš VI likewise creates man \enquote{That the service of the gods may be established}~\pid{[NE57]} \parencite{king1902seven}. The labour-substitution motif is authentic and ancient; the astronaut and biotechnological framing is Sitchin's. Heiser, himself sympathetic to ufology, could find no specialist in Sumerian or Akkadian who endorses Sitchin's work.

\subsubsection{The apkallu and the origin of the Watchers}

The strongest explanatory model for the Enochic material is literary diffusion and \emph{inversion} of Mesopotamian sage lore. \textcite{kilmer1987nephilim} first linked the Nephilim with the \emph{apkallu}, the seven antediluvian sages, whose post-flood successors were of mixed descent. \textcite{annus2010watchers} argues that the Watcher and giant mythology derives from inverted versions of apkallu myths, noting Mesopotamian traditions that already viewed the sages as dangerous. Specific convergences support the model: Gilgamesh and Humbaba appear by name in the Qumran Book of Giants (\enquote{told them what Gilgamesh said to him}~\pid{[NE33]}; \cite{reeves1992jewish,stuckenbruck1997giants}); the Astronomical Book depends on Babylonian schematic astronomy \parencite{neugebauer1981astronomical,bendov2008head,drawnel2011astronomical}; Enoch mirrors the seventh antediluvian king Enmeduranki \parencite{vanderkam1984enoch,kvanvig1988roots}; and Oannes appears as U-an, first sage, in the Seleucid Uruk List of Kings and Sages \parencite{lenzi2008uruk}, while fish-cloaked apkallu are standard protective iconography \parencite{wiggermann1992protective}. The complex later moved through Mani into Iranian and Central Asian languages, \enquote{They revealed the arts in the world, and the mysteries of heaven to the men.}~\pid{[NE30]} \parencite{henning1943giants}, and into Qur'ānic exegesis; \textcite{crone2016watchers} and \textcite{reeves2015harut} argue that the Jewish Midrash of Shemhazai and Azael is probably later than, and modelled on, the Islamic tale. Diffusion does not exhaust meaning: Genesis 6 and 1~Enoch 6--11 have also been read as political allegory, priestly polemic and theodicy \parencite{nickelsburg1977myth,suter1979fallen,hanson1977rebellion,reed2005fallen}.

\emph{Summary.} The texts contain beings from ``heaven'' who descend, mate with humans and transmit knowledge; hybrid giants whose spirits persist as demons; sages from the sea who found civilisation; human ascents to a heavenly court; and a third species of normally invisible intelligent beings. They do not contain planets of origin, technology, physical craft (outside Ezekiel's analogical throne-vision), genetic manipulation or time dilation. Their astronomy is a 364-day schematic year and their cosmography a flat earth with gates---the opposite of the out-of-period knowledge H1 would predict.

%---------------------------------------------------------------------
\subsection{South and East Asia}\label{sec:asia}
%---------------------------------------------------------------------

\subsubsection{Sources and the two groups}

The Asian workstream read Indian (Vedic, epic, Purāṇic, Jain, Buddhist and medieval technical), Japanese, Korean and Tibetan material, mainly in Victorian translations: Ganguli's \emph{Mahābhārata}, Griffith's \emph{Rāmāyaṇa} and \emph{Rig Veda}, Wilson's \emph{Viṣṇu Purāṇa}, Jacobi's \emph{Kalpa Sūtra}, Rhys Davids' \emph{Dialogues of the Buddha}, Kern's \emph{Lotus Sūtra}, Chamberlain's \emph{Kojiki} and Aston's \emph{Nihongi} \parencite{ganguli1883mahabharata,griffith1870ramayan,griffith1896rigveda,wilson1840vishnu,jacobi1884jaina,rhysdavids1899dialogues,kern1884lotus,chamberlain1882kojiki,aston1896nihongi}. Its main methodological finding is that the Asian material most often cited for extraterrestrial contact falls into two different groups: genuinely old narrative and cosmological literature, and twentieth-century products---the \emph{Vaimānika Śāstra}, the ``Dropa stones'' and the Shambhala--UFO link---that are routinely retrojected onto the first.

\subsubsection{Vimānas: vivid, and divine}

Celestial vehicles are strongly present in Indian epic, as divine or magical property. Salva, lord of Saubha, attacks Dvārakā and \enquote{rose into the sky on his car of precious metals capable of going anywhere at will!}~\pid{[asia-001]}; the car, \enquote{full two miles off, it could not, O Bharata, be seen by my troops.}~\pid{[asia-002]} Indra's car, driven by Mātali, fetches Arjuna: \enquote{the car endued with great effulgence and guided by Matali, came dividing the clouds and illuminating the firmament and filling the entire welkin with its rattle deep as the roar of mighty masses of clouds.}~\pid{[asia-004]} The flying asura city Hiraṇyapura moves in every element: \enquote{And now (the city) entered unto the earth and now it rose upwards; and at one time it went in a crooked way and at another time it submerged into water.}~\pid{[asia-007]} Rāvaṇa's Pushpaka, in Griffith's rhymed verse, \enquote{Swift moving as the master chose It flew through air or sank or rose,}~\pid{[asia-012]} and is \enquote{wrought by hands divine}~\pid{[asia-013]} \parencite{griffith1870ramayan}.

The texts describe these vehicles vividly, but they consistently place their powers in \emph{māyā} (illusion) and in boons from Brahmā or Indra: the Kalakeyas' city was created \enquote{by Brahma himself}~\pid{[asia-006]}, and it finally ``vanished, like a cloud-constructed city'', the image of the \emph{gandharva-nagara} or mirage. Vedic hymns frame the Aśvins' chariot as made by the Ṛbhus \enquote{out of your mind, by thought}~\pid{[asia-015]}, and Griffith's \enquote{animated vessels}~\pid{[asia-016]} stack ship, bird and chariot images in a way that fits Vedic poetics rather than engineering. In Jain usage \emph{vimāna} names the heavenly palaces of the gods: Mahāvīra \enquote{descended from the great Vimâna, the all-victorious and all-prosperous Pushpottara}~\pid{[asia-023]} \parencite{jacobi1884jaina}. The word's core sense is ``palace, abode''; ``vehicle'' is secondary.

The one genuinely premodern \emph{technical} description of a flying device is King Bhoja's (eleventh century) wooden bird, summarised by Raghavan: \enquote{In the bowels of the structure, below, is to be a fire-chamber with mercury placed over a flame. The power generated by the heated mercury, helped by the concurrent action of the wings which are flapped by a rider inside, makes the yantra go up and travel far}~\pid{[asia-036]} \parencite{raghavan1952yantras}. It is a human-made automaton within an Indian \emph{yantra} tradition that Sanskrit authors themselves associated with the Yavanas, i.e. Hellenistic automata: \enquote{it is the Yavanas who know the Akasa-yantras}~\pid{[asia-037]}. It has nothing to do with non-human visitors, and popular ``mercury vortex engine'' renderings add vocabulary not in Raghavan's paraphrase.

\subsubsection{The Vaimānika Śāstra}

The \emph{Vaimānika Śāstra} presents itself as \enquote{revealed to venerable Pandit SUBBARAYA SASTRY}~\pid{[asia-029]} and promises \enquote{comfortable travel in the sky from world to world}~\pid{[asia-030]} \parencite{josyer1973vymanika}. Its provenance is modern on every account. The pro-authenticity 2001 AR\&DB study reports that the verses \enquote{were recorded in 23 exercise books during the period 1903 to1918.}~\pid{[asia-031]} \parencite{rao2001vymanika}; the sacred-texts preface states that \enquote{The Vymanika Shastra was first committed to writing between 1918 and 1923, and nobody is claiming that it came from some mysterious antique manuscript.}~\pid{[asia-035]} \parencite{hare_vs_preface}. The 1974 study by engineers of the Indian Institute of Science described inspired dictation---\enquote{used to spell out shlokas (verses) whenever he got inspiration}~\pid{[asia-032]}---and concluded: 
\begin{quote}\small Thus the work "Vymanika Shastra" was brought into existence sometime between 1900 and 1922 by Pandit Subbaraya Shastry by techniques unclear to us at the moment. The only evidence in favour of Maharshi Bhardwaja being the author is the textual statement and nothing more.\par\hfill{\footnotesize --- A critical study of the work "Vymanika Shastra", Mukunda et al. 1974, sec. 1.6 Conclusions \pid{[asia-033]}}\end{quote}
 They judged the craft aeronautically unworkable and observed that the Pushpaka of the \emph{Rāmāyaṇa} \enquote{has no flying qualities except possibly by invocation of 'mantras' or 'tantras'.}~\pid{[asia-034]} \parencite{mukunda1974critical}. The text belongs with twentieth-century revelatory literature, beside Theosophy and \emph{Oahspe}, not with Vedic testimony.

\subsubsection{Heavenly descent and dynastic charters}

Ninigi's descent from Takamagahara founds the Japanese imperial line: he \enquote{pushing asunder the eight-fold heavenly spreading clouds}~\pid{[asia-040]} set off \enquote{floating shut up in the Floating Bridge of Heaven}~\pid{[asia-040]} \parencite{chamberlain1882kojiki}. That vehicle-like rendering is a translation artefact, as Chamberlain's own note admits: 
\begin{quote}\small The translator has adopted the interpretation proposed by Hirata, the only commentator who gives an acceptable view of this extremely difficult clause, which Motowori admitted that he did not understand. It must be remembered that Hirata identifies the "Floating Bridge of Heaven" with the "Heavenly Rock-Boat."\par\hfill{\footnotesize --- Kojiki, Section XXXIV, Chamberlain note 4, Kojiki Sect. XXXIV, n. 4 (p. 135) \pid{[asia-041]}}\end{quote}
 The \emph{Nihongi} has Nigi-haya-hi \enquote{riding in a Rock-boat of Heaven}~\pid{[asia-043]} who then marries a local chieftain's sister \parencite{aston1896nihongi}. In Korea, \enquote{The Old Record notes that in ancient times Hwanin's son, Hwanung, wished to descend from heaven and live in the world of human beings.}~\pid{[asia-046]} \parencite{lee1993sourcebook}; and Haemosu, father of Chumong, \enquote{He came down through the air in a five-dragon chariot, with a retinue of hundreds, robes streaming, riding on swans.}~\pid{[asia-047]} \parencite{rutt1973lay}. Across Northeast Asia the pattern ``heavenly ruler descends and founds the state'' is so common---with solar-egg births, bear totems and regalia---that diffusion within a shared sacral-kingship complex is the most economical explanation (H2).

\subsubsection{Hybrid offspring, time dilation, plurality}

Divine paternity is structural in India and Korea: the Pāṇḍavas and Karṇa are sons of gods (Kuntī \enquote{summoned the god Arka (Sun)}~\pid{[asia-009]}). The Jain embryo transfer by Harinegameṣin, a favourite of ``genetic intervention'' readings, is caste ideology in context---a tīrthaṅkara must be born a Kṣatriya---and Digambara Jains reject the episode.

The Raivata/Kakudmin story of the \emph{Viṣṇu Purāṇa} is a clear case of time dilation: 
\begin{quote}\small When he arrived, the quiristers Háhá, Húhú, and others, were singing before Brahmá; and Raivata, waiting till they had finished, imagined the ages that elapsed during their performance to be but as a moment.\par\hfill{\footnotesize --- Vishnu Purana, Book IV, ch. 1, VP 4.1 (Wilson p. 353ff.) \pid{[asia-017]}}\end{quote}
 He returns to find \enquote{the race of men dwindled in stature, reduced in vigour, and enfeebled in intellect.}~\pid{[asia-018]} Wilson already explained the device: \enquote{The object of this legend, which is told by most of the authorities, is obviously to account for the anachronism of making Balaráma cotemporary with Raivata}~\pid{[asia-019]} \parencite{wilson1840vishnu}. In Japan, by contrast, the \emph{Nihongi}'s Urashima entry has the otherworld voyage to Hōrai but no time lapse; the three-hundred-year motif appears only in the Tango \emph{fudoki} fragment and the \emph{Man'yōshū} \parencite{nicolae2017myth}.

The strongest and most original Asian contribution is plurality of worlds, but it is metaphysical. The \emph{Viṣṇu Purāṇa} speaks of cosmic eggs \enquote{of which there are thousands and tens of thousands, and millions and thousands of millions}~\pid{[asia-022]}; the \emph{Lotus Sūtra} of \enquote{eighteen hundred thousand Buddha-fields}~\pid{[asia-027]} and of bodhisattvas \enquote{who had flocked from other worlds}~\pid{[asia-028]} \parencite{kern1884lotus}; the Pali Brahmajāla of luminous devas who \enquote{dwell made of mind, feeding on joy, radiating light from themselves, traversing the air}~\pid{[asia-025]} \parencite{rhysdavids1899dialogues}. Arjuna, riding Indra's car, sees that the stars are vast and distant: \enquote{And those brilliant regions that are seen from the earth in the form of stars, like lamps (in the sky)--so small in consequence of their distance, though very large--were beheld by the son of Pandu}~\pid{[asia-005]} But Mātali explains that they are the stations of meritorious persons. These are cosmologies of rebirth and liberation; their inhabitants are gods, buddhas and bodhisattvas, and contact is visionary or meditative.

\subsubsection{Modern constructions: ``nuclear war'', Shambhala, Dropa}

The famous claim that the \emph{Mahābhārata} describes survivors of a nuclear weapon losing hair and nails rests on an omen: \enquote{At night, the rats and mice ate away the hair and nails of slumbering men.}~\pid{[asia-011]} The widely quoted \enquote{single projectile charged with all the power of the universe}~\pid{[asia-050]} does not occur in Ganguli at all; it is a modern paraphrase circulated by Childress and others \parencite{childress_aircraft}. Berzin traces the Shambhala--UFO association to a chain of modern reinterpretation (Blavatsky, Roerich, Ossendowski, Saint-Yves d'Alveydre, a Brazilian Theosophist's hollow-earth saucers) whose only textual hook is an allegorical apocalyptic war in the \emph{Vimalaprabhā} commentary; the Third Panchen Lama's guide, he notes, makes the journey \enquote{actually an inner quest}~\pid{[asia-049]} \parencite{berzin2003mistaken}. The ``Dropa stones''---\enquote{a space probe from a distant planet that crash-landed in the Baian-Kara-Ula mountains}~\pid{[asia-051]} \parencite{hausdorf1998dropa}---first appear in a 1962 German vegetarian magazine and have no Chinese textual or archaeological basis \parencite{creighton1973book,fitzpatrickmatthews_dropa}.

\emph{Summary.} Present in the texts: vivid divine vehicles and flying cities within a logic of boon and \emph{māyā}; heavenly descent as dynastic charter; divine paternity; one Purāṇic time-dilation tale; rich metaphysical plurality of worlds; a pervasive taxonomy of non-human beings (nāgas, gandharvas, apsarases, asuras). Modern: the \emph{Vaimānika Śāstra}, ``ancient nuclear war'', Saubha as mothership \parencite{thompson1993alien}, Shambhala saucers and the Dropa discs. No passage requires literal contact.

%---------------------------------------------------------------------
\subsection{Celtic and Northern Europe}\label{sec:celtic}
%---------------------------------------------------------------------

\subsubsection{Three bodies of material}

NHI-oriented literature routinely merges three different bodies of material: medieval Irish learned literature about the Túatha Dé Danann and the \emph{síd} peoples (\emph{Lebor Gabála Érenn}, \emph{Cath Maige Tuired}, \emph{Immram Brain}, \emph{Echtrae Chonnlai}); early-modern and modern vernacular fairy belief in Gaelic Scotland, Ireland, Man, Wales, Cornwall and Brittany, witnessed above all by Kirk's \emph{Secret Common-Wealth} (1691/92), the Child ballads, and Evans-Wentz's field interviews of 1907--10; and the modern comparative thesis, running from the Paracelsian \emph{Comte de Gabalis} (1670) through Evans-Wentz to Vallée, that fairy lore and abduction reports record one phenomenon \parencite{macalister1941lge,gray1982cmt,meyer1895,kirk1893,child1882,evanswentz1911,villars1670,vallee1969}. The first is monastic literature; \emph{Lebor Gabála} is an eleventh-century learned synthesis fitting Irish origins into biblical world history \parencite{carey1994}, and the name \emph{Túatha Dé Danann} itself dates only from about the end of the twelfth century \parencite{williams2016}. Lady Gregory's \emph{Gods and Fighting Men} (1904), through which most readers meet this material, is a harmonising retelling, not a translation \parencite{gregory1904}.

\subsubsection{Hidden peoples: the core idea}

The strongest and most pervasive Celtic motif is a co-present, normally invisible people living in mounds, lakes and islands (\textsc{dim}). After the Milesian conquest, in Gregory's retelling, Manannán hides the Túatha Dé: 
\begin{quote}\small So he chose out the most beautiful of the hills and valleys of Ireland for them to settle in; and he put hidden walls about them, that no man could see through, but they themselves could see through them and pass through them.\par\hfill{\footnotesize --- Gods and Fighting Men, Part I Bk IV 'Bodb Dearg', Part I, Book IV, ch. 1 (opening) \pid{[celtic-14]}}\end{quote}
 Kirk's \emph{sìth} 
\begin{quote}\small are said to be of a midle Nature betuixt Man and Angel, as were Dæmons thought to be of old; of intelligent fluidious (?) Spirits, and light changable Bodies, (lyke those called Astral,) somewhat of the Nature of a condensed Cloud, and best seen in Twilight.\par\hfill{\footnotesize --- The Secret Common-Wealth, ch. 1 'Of the Subterranean Inhabitants', ch. 1 (Lang ed. pp. 5-7) \pid{[celtic-22]}}\end{quote}
 They migrate at the quarter-days, and \enquote{Their chamælion-lyke Bodies swim in the Air near the Earth with Bag and Bagadge}~\pid{[celtic-23]} \parencite{kirk1893}. Eddic and Snorrian elves offer a Northern parallel: light-elves in Álfheimr, while \enquote{the Dark-Elves dwell down in the earth}~\pid{[celtic-44]} \parencite{brodeur1916}. This is a cosmology of adjacent realms, not of other planets.

\subsubsection{Abduction and the otherworld journey}

Abduction is strong, in two forms. The literary \emph{echtrae} is a summons: \enquote{the woman from unknown lands sang on the floor of the house to Bran son of Febal}~\pid{[celtic-17]}, appearing inside a closed fortress \enquote{since the ramparts were closed.}~\pid{[celtic-17]} \parencite{meyer1895}; Connla's visitor comes \enquote{from Tir-nam-Beo, the Land of the Ever-Living Ones, where no death comes.}~\pid{[celtic-15]} The folk belief is a \emph{taking}. Kirk reports: 
\begin{quote}\small Women are yet alive who tell they were taken away when in Child-bed to nurse Fairie Children, a lingering voracious Image of their (them?) being left in their place, (like their Reflexion in a Mirrour,)\par\hfill{\footnotesize --- The Secret Common-Wealth, ch. 4, ch. 4 (Lang ed. pp. 12-13) \pid{[celtic-24]}}\end{quote}
 The ballads give it verse form: \enquote{I am but the queen of fair Elfland, And I'm come here for to visit thee.}~\pid{[celtic-28]}; \enquote{The Queen o Fairies she caught me, In yon green hill to dwell.}~\pid{[celtic-29]} \parencite{child1882}. An Irish seer told Evans-Wentz that the fairies \enquote{take the whole body and soul of young and intellectual people who are interesting}~\pid{[celtic-31]}, and that \enquote{Once they take you and you taste food in their palace you cannot come back.}~\pid{[celtic-31]} \parencite{evanswentz1911}. Kirk's account of a woman \enquote{returning to her Husband after two Years Space}~\pid{[celtic-26]}, with the fairy-ointment motif and a hall lit without lamp, is the closest early-modern analogue to a ``returned abductee'' narrative.

\subsubsection{The supernatural lapse of time}

Time dilation is strong, early and international. \emph{Immram Brain}, dated by Meyer to the seventh century and by Mac Mathúna to the early eighth \parencite{meyer1895,macmathuna1985}, ends with the returning voyagers unknown: 
\begin{quote}\small Said Bran: 'I am Bran the son of Febal,' saith he. However, the other saith: 'We do not know such a one. though the Voyage of Bran is in our ancient stories.' ... As soon as he touched the earth of Ireland, forthwith he was a heap of ashes\par\hfill{\footnotesize --- Immram Brain, §§64-65 \pid{[celtic-19]}}\end{quote}
 Oisín, in Gregory's retelling of the eighteenth-century lay, touches the ground: \enquote{And in the minute all the years came on me}~\pid{[celtic-21]}. The motif continues into twentieth-century testimony. From Scotland: \enquote{when he returned two generations had disappeared during the lapse of time}~\pid{[celtic-36]}, though \enquote{to him the time had seemed only a few hours}~\pid{[celtic-36]}; from Wales: \enquote{He told his mother how he had been away for a fortnight, as he thought, but she told him it had been for two years.}~\pid{[celtic-38]} \parencite{evanswentz1911}. Hartland's chapter on the supernatural lapse of time, with Walter Map's King Herla (c.~1181--93), shows the motif is not specifically Celtic: \enquote{The king was stupefied, for he deemed he had only been away three days}~\pid{[celtic-41]} \parencite{hartland1891,map1983}.

\subsubsection{Unions, offspring and changelings}

\textsc{hyb} is strong but covers several different things: divine or otherworld paternity of kings and heroes (Manannán and Mongán, Elatha and Bres, Ríg fathering the social classes); fairy lovers and succubi; and changelings. A changeling is a \emph{substitution}, not a hybrid, and divine paternity is a sovereignty charter. Hartland's conclusion on changelings remains the strongest naturalistic counter to ``hybrid'' readings: \enquote{Making allowance for the influence of imagination, there can be no doubt, on comparison of these passages, that the children to whom the character of changelings was ascribed were invariably deformed or diseased.}~\pid{[celtic-42]} \parencite{hartland1891}; later work reached similar conclusions through disease, disability and violence against children \parencite{eberly1988,bourke1999}.

\subsubsection{Sky and vehicles: weak}

The Túatha Dé arrive, according to \emph{Lebor Gabála}, \enquote{in dark clouds}~\pid{[celtic-01]}, and \enquote{brought a darkness over the sun for three days and three nights.}~\pid{[celtic-01]} But \emph{Cath Maige Tuired} itself rationalises the clouds: 
\begin{quote}\small they at once burned their boats so that they would not think of fleeing to them. The smoke and the mist which came from the ships filled the land and the air which was near them. For that reason it has been thought that they arrived in clouds of mist.\par\hfill{\footnotesize --- Cath Maige Tuired, §9 \pid{[celtic-05]}}\end{quote}
 \parencite{gray1982cmt}. A \emph{Lebor Gabála} poem admits that \enquote{the truth was not known beneath the sky of stars, whether they were of heaven or of earth.}~\pid{[celtic-02]} Evans-Wentz cites the Book of the Dun Cow to the effect that the learned thought they \enquote{came from heaven, on account of their intelligence and for the excellence of their knowledge}~\pid{[celtic-03]}. This is a debate in Christian angelology about the status of pre-Christian gods, not evidence of aerial arrival. Gregory's \enquote{came through the air and the high air to Ireland.}~\pid{[celtic-12]} heightens the aerial element in retelling. The otherworld vessels (Elatha's silver ship, Connla's shining boat, Manannán's chariot) travel on the sea. The one genuine ``ships in the clouds'' belief in the corpus comes from Carolingian Francia, not from Celtic sources (Section~\ref{sec:magonia}).

\subsubsection{Culture-bringers, war, planets and channelling}

Lug, the \emph{Samildánach}, is \enquote{the man of each and every art.}~\pid{[celtic-09]}, matching the Gaulish Lugus whom Caesar called the inventor of all arts---inherited Celtic theology \parencite{caesar_bg}. The Túatha Dé themselves learned their arts \enquote{in the northern islands of the world, studying occult lore and sorcery}~\pid{[celtic-06]}: they are learners and possessors of knowledge, not sky-teachers of humanity. The war of the Túatha Dé and the Fomoire is an intra-supernatural conflict of Indo-European type; the ``beam-weapon'' reading of Balor's eye is modern. Plurality of worlds and channelling are essentially absent from tradition. The only ``planets'' statement is one Sligo informant's private opinion, c.~1909: 
\begin{quote}\small The gentry are the most noble tribe of all; and they are a big race who came from the planets--according to my idea; they usually appear white.\par\hfill{\footnotesize --- The Fairy-Faith in Celtic Countries, ch. II: testimony of Patrick Wate, p. 53 \pid{[celtic-32]}}\end{quote}
 He gave it alongside a traditional fallen-angel cosmology, and it most likely reflects the popular astronomy of the period. \textsc{chan} appears only in the visions of the Theosophist AE (George Russell), whose beings have \enquote{a collective life, so unindividualized and so calm}~\pid{[celtic-34]}---a description that anticipates late-twentieth-century ``Grey'' lore but whose source is Theosophical monism.

\subsubsection{The Magonia chain}\label{sec:magonia}

Two transmission chains are especially instructive. The first runs from Agobard of Lyon to Villars to Vallée. Agobard (c.~815) reports only a popular belief and a mob: 
\begin{quote}\small they believe and say that there is a certain region, which is called Magonia, from which ships come in the clouds. In these ships the crops that fell because of hail and were lost in storms are carried back into that region; evidently these aerial sailors make a payment to the storm-makers\par\hfill{\footnotesize --- Agobard of Lyon, De grandine et tonitruis, ch. 2 \pid{[celtic-46]}}\end{quote}
 \parencite{agobard_lewis2001}. Villars (1670) turns this into sylphs \enquote{on wonderfully constructed aerial ships}~\pid{[celtic-47]}, and a Lyons episode in which people \enquote{had been carried away a short time since by miraculous men who had shown them unheard-of marvels}~\pid{[celtic-48]} \parencite{villars1670}. The abduction narrative is an esoteric addition of 1670, and Vallée's title inherits it.

The second runs from Evans-Wentz to Vallée. Evans-Wentz's footnote on a Barra fairy-mistress tale is cautious: \enquote{This curious tale suggests that certain of the fairy women who entice mortals to their love in modern times are much the same, if not the same, as the succubi of Middle-Age mystics.}~\pid{[celtic-37]} The online reproduction of \emph{Passport to Magonia} renders this as ``proves conclusively'', and the informants' names differ. Whether the discrepancy lies in Vallée's 1969 text or in the online transcription could not be determined, and the quotation must be checked against the Regnery printing before it is used as evidence of Vallée's practice. Evans-Wentz himself linked folk belief to what would now be called NHI theory. His conclusion holds that \enquote{Fairies exist, because in all essentials they appear to be the same as the intelligent forces now recognized by psychical researchers}~\pid{[celtic-40]} \parencite{evanswentz1911}. Vallée's interdimensional ``control system'' descends directly from this argument; he could write that \enquote{There is no gap between the fairy-faith and ufology regarding the sexual question.}~\pid{[celtic-51]} \parencite{vallee1969}, and later that \enquote{It is difficult to find a culture on earth that does not have an ancient tradition of little people that fly through the sky and abduct humans}~\pid{[celtic-52]} \parencite{vallee1990}. The latter generalisation is not borne out by the present corpus: flight is precisely what Celtic fairies rarely do.

\emph{Summary.} Genuinely in the tradition: a co-present hidden people; abduction of new mothers, infants, brides and gifted youths; lapse of time and contact taboos (food, dismounting, touching the earth); sexual liaisons and supernatural paternity; luminous beings and sourceless light; elf-shot ``physical traces'' that are in fact Neolithic arrowheads (celtic-27); and medieval uncertainty about whether the Túatha Dé were angels, demons or humans. Modern: Túatha Dé spacecraft, the Otherworld as another planet, changelings as hybrids, and relativistic time dilation.

%---------------------------------------------------------------------
\subsection{The Americas and the Pacific}\label{sec:americas}
%---------------------------------------------------------------------

\subsubsection{Sources}

This area spans Indigenous North America (Ojibwa, Arapaho, Haudenosaunee, Cherokee, Hopi, Diné, Acjachemen), K'iche' and Yucatec Maya, Nahua, Inca, Kayapó, Māori and Aboriginal Australia. Almost all ``primary'' material is mediated: colonial-era alphabetic texts written after the Conquest (the \emph{Popol Vuh}, Landa, Sarmiento), or ethnographic texts produced by outsiders between about 1820 and 1980, usually from a few named informants \parencite{hewitt1903iroquoian,mooney1900myths,voth1905traditions,boscana1846chinigchinich,smith1913lore,parker1905euahlayi,metraux1960mythes}. Of its 50 passages, 29 are ethnographic and 8 primary. Its motif profile is distinctive: \textsc{sky} and \textsc{cult} 18 each, \textsc{dim} 14, \textsc{star} 12, and no \textsc{time} at all; all five \textsc{veh} tags fall on fringe or late sources.

\subsubsection{Layered worlds and sky-people}

Most North American cosmologies describe a layered universe whose sky is a land with its own people. The Iroquoian sky-world is formulaic: \enquote{Man-beings dwell in the sky, on the farther side of the visible sky [the ground separating this from the world above it].}~\pid{[AP05]} \parencite{hewitt1903iroquoian}; the Seneca text opens: \enquote{A long time ago human beings lived high up in what is now called heaven. They had a great and illustrious chief.}~\pid{[AP04]} \parencite{thompson1929tales}. Sky Woman falls through a hole towards a watery world: \enquote{there seemed to be a lake at the spot toward which she was falling. There was nowhere any earth.}~\pid{[AP06]} Hewitt's own philological note is crucial: Onondaga \emph{oñgwe} (``man-being'') covers personified powers such as Daylight, Wind and Earthquake, and he describes some as \enquote{Some of these beings were mere fictions, figures of speech made concrete and objective.}~\pid{[AP07]} English ``sky people'' therefore overstates how humanoid these beings are---a translation artefact that NHI readings exploit. The Kayapó offer a parallel. Métraux's text opens \enquote{Autrefois, tous les hommes vivaient au ciel.}~\pid{[AP34]} (``Once all people lived in the sky''); an old man digs through the sky floor, and people descend by rope, while those who remain are the stars \parencite{metraux1960mythes}. The sky-hole-and-rope sequence also appears in the Star Husband tale and in Sky Woman; Thompson noted that \enquote{a star-world, a sky window, a rope to the sky, a rainbow-bridge to the upper world are to be found everywhere.}~\pid{[AP03]} Across the hemisphere this is best read as a deep, shared layered-world cosmology (H2/H3).

\subsubsection{Star-persons and star marriage}

The Star Husband tale is the best-documented sky-marriage narrative in the Americas. In the Timagami Ojibwa version, two girls wish for star husbands: 
\begin{quote}\small When they awoke, they found themselves in another world, the star world. There were four of them there, the two girls and the two stars who had become men.\par\hfill{\footnotesize --- Stith Thompson, Ch. V, No. L, pp. 126-128 \pid{[AP01]}}\end{quote}
 Thompson's historic-geographic monograph compiled 86 versions from 44 tribes in nine culture areas and argued for diffusion from the central Plains \parencite{thompson1953star}; \textcite{dundes1965study} and \textcite{levistrauss1968origine} reanalysed the corpus structurally. The idiom is kinship and moral pedagogy, not technology. Cherokee informants disagreed about what stars are: \enquote{There are different opinions about the stars. Some say they are balls of light, others say they are human, but most people say they are living creatures covered with luminous fur or feathers.}~\pid{[AP08]} \parencite{mooney1900myths}. Pleiades myths in which youths rise into the sky recur among the Cherokee (\enquote{with every round they rose higher and higher in the air.}~\pid{[AP09]}), the Maya (\enquote{the four hundred boys died, and it is said that they became the group of stars}~\pid{[AP26]}) and the Euahlayi (\enquote{The Pleiades are seven sisters}~\pid{[AP46]}; \cite{parker1905euahlayi}). In each case humans become stars; stars are not the origin of humans. Fringe ``Pleiadian ancestry'' readings reverse the direction.

\subsubsection{Culture-bringers, including the one from the stars}

The most explicit star-origin teacher in the entire corpus is the Acjachemen Chinigchinich, recorded by the Franciscan Gerónimo Boscana c.~1822--25: 
\begin{quote}\small He said that he had come from the stars to teach them those things of which they were ignorant.\par\hfill{\footnotesize --- Gerónimo Boscana, Chinigchinich, Ch. III (Creation according to those on the sea-coast), p. 1 \pid{[AP21]}}\end{quote}
 Elsewhere he says, \enquote{My habitation is above.}~\pid{[AP22]} \parencite{boscana1846chinigchinich}. This is the one traditional passage in the study in which a teacher explicitly comes ``from the stars''. It calls for caution on two counts: Boscana's Franciscan framing, and the possibility that the Chinigchinich cult was an early-historic revitalisation movement \parencite{kroeber1925handbook,bean1978cults}. In a cosmology in which stars are persons and ``above'' is a populated tier, it is also naturally read as a sky-deity's self-description.

Mesoamerican and Andean ``visitors'' are largely colonial constructs. Landa's Kukulcán is a human lord: \enquote{They say that he came from the West}~\pid{[AP29]}, founded Mayapán and returned to Mexico with no promise of return \parencite{landa1937yucatan}. The ``fair, bearded'' Quetzalcoatl derives from Franciscan-era writers such as Torquemada (1615), and the ``returning white god'' story took shape after the Conquest \parencite{townsend2003burying}. Sarmiento's Viracocha, \enquote{white and dressed in a white robe like an alb secured round the waist}~\pid{[AP32]} and carrying \enquote{a staff and a book in his hands.}~\pid{[AP32]}, wears obviously Christian-apostolic imagery, as does his prophecy of false Viracochas and future teaching messengers \parencite{sarmiento1907history}. The \emph{Popol Vuh}, written c.~1554--58, declares itself composed \enquote{under the Law of God and Christianity}~\pid{[AP23]} \parencite{goetz1950popol}.

\subsubsection{A limit on knowledge}

The \emph{Popol Vuh} contains the clearest divine limitation of knowledge in the Americas. The first maize-people \enquote{saw and instantly they could see far, they succeeded in seeing, they succeeded in knowing all that there is in the world.}~\pid{[AP24]} Then: \enquote{Then the Heart of Heaven blew mist into their eyes, which clouded their sight as. when a mirror is breathed upon. Their eyes were covered and they could see only what was close, only that was clear to them.}~\pid{[AP25]} The structure resembles Genesis 3:22, which raises the question of colonial influence; and it inverts the Watchers/Prometheus pattern---knowledge is withdrawn by the gods rather than leaked or stolen.

\subsubsection{Hidden peoples and emergence}

The Cherokee Nûñnë'hï were \enquote{invisible excepting when they p. 331 wanted to be seen}~\pid{[AP10]} and then looked just like other Indians; the Little People are \enquote{little fellows, hardly reaching up to a man's knee}~\pid{[AP11]} \parencite{mooney1900myths}. The resemblance to Irish \emph{síd} lore is striking and is usually explained by similar ways of thinking about liminal landscape, possibly with some Scots-Irish contact influence. Hopi and Diné emergence narratives describe successive lower worlds: \enquote{A long time ago the people were living below.}~\pid{[AP13]} \parencite{voth1905traditions}; this is \enquote{the fourth plane on which mankind has existed.}~\pid{[AP15]} \parencite{lockett1933unwritten}. These are strata inside the earth, not other planets. Mooney's ``Star Feathers'' supplies the area's only channelling passage, and it is a parody: a man \enquote{asserted that he had been up to the sky and that these were star feathers}~\pid{[AP12]} and delivered a message \enquote{which he pretended was a message he had received from the star spirits}~\pid{[AP12]}, until he was exposed. The tradition itself contained critical discourse about contactee-style claims.

\subsubsection{Four fringe claims traced to their sources}

\emph{The Hopi Ant People.} Neither Voth nor Lockett mentions Ant People sheltering humans; Lockett describes kachinas as \enquote{Kachinas, spirits of ancestors and some other beings, with powers good and bad.}~\pid{[AP14]} The Ant People and flesh-and-blood kachinas from a planet---\enquote{They came from a distant planet the Hopi called Tóonáotakha, located outside our solar system.}~\pid{[AP16]} \parencite{lloyd2017watchers}---go back to Oswald ``White Bear'' Fredericks, informant for Waters' \emph{Book of the Hopi} (1963) and Blumrich's \emph{Kásskara} (1979); Geertz shows that White Bear was not initiated into kiva rites \parencite{waters1963book,blumrich1979kasskara,geertz1983book,geertz1994invention}. The ``anu = Anunnaki'' etymology---\enquote{may be related to Sitchin's Anunnaki.}~\pid{[AP17]} \parencite{david2002anthills}---is a chance resemblance. In Diné tradition ``ant people'' are insect-persons of the lower worlds (AP19).

\emph{Kayapó Bep Kororoti.} Métraux's 1960 field text is an origin of rain and lightning: Bepkororoti, cheated of tapir meat, \enquote{Bepkororoti protesta et se fâcha, mais rien n'y fit.}~\pid{[AP35]} calls down storm, and ascends to the sky. \textcite{posey1983keeping} independently records a living ritual: honey is left \enquote{for Bep-kororoti, a powerful shaman who was taken into the sky in a flash of lightning.}~\pid{[AP36]} The spacesuit and teaching features first appear in print in von Däniken (1973) \parencite{vondaniken1973welt}, whose source was the indigenist João Américo Peret; Peret's own article followed in 1977, glossing the straw costume as ``overall and helmet'' and adding that \enquote{uma especie de estrela pousava}~\pid{[AP38]} (``a kind of star landed'') \parencite{peret1977bep}. The version reproduced online contains an editorial interjection inside supposedly indigenous speech: \enquote{The warrior who had come from the cosmos (or the future? - Red.) must have laughed at the weakness of those who fought against him.}~\pid{[AP40]} That is decisive evidence of contamination.

\emph{Wandjina.} The ethnographic doctrine is that \enquote{the Wandjina put their own images on the cave walls before they returned to the spirit world}~\pid{[AP47]}, and repainting them sustains life \parencite{frederick2009wandjina,crawford1968art}. The astronaut reading grows out of colonial denial of Indigenous authorship (Grey in 1841 thought it improbable that a ``self-taught savage'' had painted them) and was taken up by von Däniken: \enquote{The characteristic semicircular headdress is often referred to as an extraterrestrial’s helmet.}~\pid{[AP48]} \parencite{grey1841journals,frederick2009wandjina}. An anonymous Spanish page turns the self-painting doctrine into \enquote{unos seres que descendieron a la Tierra en tiempos remotos.}~\pid{[AP49]} \parencite{anon_wandjinas}.

\emph{Quetzalcoatl as extraterrestrial.} The argument jumps from Torquemada's \enquote{a fair and ruddy complexioned man with a long beard.}~\pid{[AP31]} to modern UAP news, with no indigenous evidence in between \parencite{walia2018quetzalcoatl}.

\subsubsection{Source-critical cautions}

The Whare-wānanga corpus, in which Tāne ascends to Io for the baskets of knowledge (AP41--42) and each heaven has \enquote{its own form of life suited to each.}~\pid{[AP44]}, is the area's nearest approach to a plurality of inhabited worlds---and its least secure passage: Simmons and Biggs showed that Whatahoro assembled the manuscripts from various sources \parencite{smith1913lore,simmons1970sources}. The south-eastern Australian All-Father was reshaped in response to invasion and missions, although attestations of ``Piame'' (1832) predate the major missions \parencite{swain1993place,henderson1832observations}. Any ``bearded white visitor'' in colonial Mesoamerican or Andean texts should be treated as a colonial layer unless pre-Hispanic iconography supports it.

\emph{Summary.} Present in the texts: layered worlds with sky-people, star-persons who marry humans, humans who become stars, culture-bringers from ``above'' (one explicitly ``from the stars''), divine limitation of knowledge, and invisible hidden peoples. No text in the early ethnographic or colonial layer describes a vehicle, a technology, a planet of origin, or time dilation. Nearly every explicitly extraterrestrial element (White Bear/Waters/Blumrich, von Däniken/Peret, Wandjina helmets) appears after 1960.

%=====================================================================
\section{Cross-cultural motif analysis}\label{sec:crosscultural}
%=====================================================================

\subsection{The motif matrix}

Table~\ref{tab:matrix} and Figure~\ref{fig:heat} give the distribution of motif tags by area. Because a passage can carry several motifs, the cells count tags, not passages. Panel (a) of Figure~\ref{fig:heat} counts all primary and ethnographic (P/E) passages; panel (b) counts only those rated strong. Two caveats apply. First, the counts reflect the workstreams' sampling of the most relevant passages, not the frequency of motifs in the traditions as wholes; they measure the \emph{shape} of the evidence, not its prevalence. Second, the P/E filter under-represents Africa, because the richest African sky-people material (Werner's Chaga and Rwandan texts) is typed S as a secondary compilation; Africa's strong P/E row contains only 13 tags, whereas all African passages together carry 73.

\begin{table}[ht]
\centering\small
\caption{Motif tags by area. Each cell gives the number of tags on all verified passages and, in brackets, on strong primary/ethnographic (P/E) passages only.}\label{tab:matrix}
\setlength{\tabcolsep}{3.2pt}
\resizebox{\textwidth}{!}{%
\begin{tabular}{@{}lccccccccccccc@{}}
\toprule
Area & \textsc{sky} & \textsc{cult} & \textsc{forb} & \textsc{hyb} & \textsc{veh} & \textsc{star} & \textsc{time} & \textsc{abd} & \textsc{plur} & \textsc{dim} & \textsc{chan} & \textsc{war} & Total \\
\midrule
Near East & 14 (10) & 14 (8) & 12 (9) & 17 (13) & 9 (3) & 4 (3) & -- & 14 (6) & 2 (0) & 9 (3) & 1 (1) & 12 (9) & 108 (65) \\
Africa & 14 (4) & 10 (3) & -- & 6 (1) & 5 (2) & 12 (1) & -- & 10 (1) & 3 (0) & 7 (0) & 3 (0) & 3 (1) & 73 (13) \\
Asia & 15 (8) & 4 (2) & -- & 6 (5) & 27 (13) & 3 (0) & 3 (2) & 9 (3) & 7 (3) & 8 (1) & 6 (2) & 9 (6) & 97 (45) \\
Celtic & 6 (2) & 6 (2) & 1 (0) & 14 (5) & 8 (3) & -- & 7 (5) & 19 (11) & 1 (0) & 23 (8) & 2 (0) & 1 (0) & 88 (36) \\
Amer./Pac. & 18 (8) & 18 (3) & 2 (2) & 2 (2) & 5 (0) & 12 (8) & -- & 7 (4) & 4 (3) & 14 (8) & 1 (1) & 1 (0) & 84 (39) \\
Modern & 25 (2) & 15 (1) & 5 (0) & 10 (0) & 20 (1) & 6 (1) & -- & 1 (0) & 19 (1) & 4 (0) & 15 (1) & 8 (1) & 128 (8) \\
\midrule
All & 92 (34) & 67 (19) & 20 (11) & 55 (26) & 74 (22) & 37 (13) & 10 (7) & 60 (25) & 36 (7) & 65 (20) & 28 (5) & 34 (17) & 578 (206) \\
\bottomrule
\end{tabular}}
\end{table}

\begin{figure}[ht]
\centering
\includegraphics[width=\textwidth]{figures/fig1_heatmap_primary.pdf}
\caption{Motif by area for primary and ethnographic evidence in the five traditional areas: (a) all P/E passages; (b) P/E passages rated strong. Cells count motif tags; ``·'' marks zero.}\label{fig:heat}
\end{figure}

\subsection{What is strong, what is absent, and where}

\paragraph{Universal or near-universal.} Sky descent (\textsc{sky}), culture-bringers (\textsc{cult}), divine--human unions (\textsc{hyb}), otherworld journeys (\textsc{abd}) and hidden peoples (\textsc{dim}) occur in all five traditional areas. They are also the motifs most obviously tied to universal social functions: legitimising rulers (heavenly descent and divine paternity as dynastic charters in Japan, Korea, Yoruba Ife and Ireland), explaining the origin of arts and crops, and accounting for death, illness, infertility and disability (the fairy taking, the changeling).

\paragraph{Forbidden knowledge is essentially Near Eastern.} Of the 14 \textsc{forb} tags on traditional P/E passages, 11 are Near Eastern; the strong cases outside the Near East are the two \emph{Popol Vuh} passages in which the gods \emph{withdraw} knowledge from humans---an inversion of, not a parallel to, the Watchers. The Ethiopic Watchers of the \emph{Kebra Nagast} teach nothing; the Túatha Dé learn occult arts rather than leak them. The specific pattern of heavenly beings illicitly transmitting civilisation's secrets is a single literary tradition (Mesopotamian sage lore inverted in Second Temple Judaism, then transmitted to Manichaeism and Islam), not a cross-cultural universal.

\paragraph{Time dilation: Celtic, Indian, and late Japanese.} \textsc{time} occurs in eight traditional P/E passages: six Celtic and two from the \emph{Viṣṇu Purāṇa}'s Raivata story. In Japan the Urashima story acquires its three-hundred-year lapse only in the \emph{fudoki} and \emph{Man'yōshū} versions; the \emph{Nihongi} entry has none. 2~Enoch's sixty days in heaven equal sixty days on earth. Africa, the Americas and the Near East supply no case. Time dilation is thus a strong but regionally clustered motif, attested early (Ireland, 7th--8th century; the Purāṇa, first millennium CE) in cultures that also conceive the otherworld as the land of the ever-living.

\paragraph{No traditional planet of origin.} No traditional text in any area names a planet or a particular star system as the origin of visiting beings. ``Heaven'' in these texts is a cosmological tier---a roof, a court, a layered sky---not a location in space. Stars appear as persons (Cherokee, /Xam, Ojibwa), as angels (1~Enoch), or as transformed humans (Pleiades myths on three continents); humans become stars far more often than stars become humans. The nearest exceptions are instructive precisely because they are so few and so qualified: Chinigchinich's claim to have come ``from the stars'' (a Franciscan-recorded, possibly early-historic cult), a Sligo informant's avowedly personal opinion that the fairies came ``from the planets'' (c.~1909), and Kaguya-hime's lunar origin in a 1911 children's adaptation. Planetary origins (Venus, Mars, Sirius, Nibiru, ``Tóonáotakha'') enter only with modern texts.

\paragraph{Vehicles are divine, magical or human-made.} Every vehicle in the primary texts is framed by divine favour, boon, mantra, \emph{māyā} or heroic property: the \emph{Kebra Nagast}'s wagons raised by an archangel, Indra's horse-drawn car, the Pushpaka built by the divine architect, flying asura cities that vanish like mirages, Japanese heavenly rock-boats, Haemosu's dragon chariot, Manannán's chariot on the sea. Ezekiel's wheels are explicitly ``likeness''. The only premodern technical description, Bhoja's mercury-heated wooden bird, is a human automaton. The two strong ``technical'' vehicle passages in the Asian P/E row are from the \emph{Vaimānika Śāstra}, a twentieth-century text. The Americas contribute no vehicle from any early source.

\paragraph{Channelling is modern.} Of the six \textsc{chan} tags on traditional P/E passages, two are the modern \emph{Vaimānika Śāstra}, two are the Theosophist AE's visions recorded in 1911, one is the Qur'ānic jinn who overhear a revelation (a reversal of the contactee pattern), and one is the Cherokee parody of a man faking messages from star spirits.

\subsection{The shape of fringe evidence}

Figure~\ref{fig:fringe} contrasts the full corpus with the fringe (F) passages alone, and Figure~\ref{fig:motiftotals} splits each motif's tags by evidential class. The contrast is sharp. Among the 187 traditional P/E passages the commonest motifs are sky descent (28\%), hidden peoples (24\%) and otherworld journeys (23\%); vehicles account for 17\% and plurality of worlds for 7\%. Among the 72 fringe passages, vehicles (42\%) and sky descent (40\%) dominate, followed by culture-bringers (28\%), plurality of worlds and hybrids (25\% each) and channelling (17\%); otherworld journeys fall to 8\% and time dilation to zero. The modern reading, in other words, does not merely reinterpret the traditional motif profile; it \emph{reweights} it, promoting vehicles, planets and messages and demoting the hidden people, the taking and the lapse of time that dominate the traditional record.

\begin{figure}[ht]
\centering
\includegraphics[width=\textwidth]{figures/fig2_heatmap_all_fringe.pdf}
\caption{(a) Motif tags for all 314 verified passages, including the modern-reception area; (b) fringe/esoteric (F) passages only.}\label{fig:fringe}
\end{figure}

\begin{figure}[ht]
\centering
\includegraphics[width=0.82\textwidth]{figures/fig4_motif_totals.pdf}
\caption{Motif tags across the corpus by evidential class: primary/ethnographic passages in the five traditional areas; secondary and critical scholarship (plus the three primary control passages in the modern area); and fringe/esoteric passages.}\label{fig:motiftotals}
\end{figure}

%=====================================================================
\section{Reception history: how the ``ancient NHI'' reading was built}\label{sec:reception}
%=====================================================================

\subsection{A datable genealogy}

The modern-reception workstream asked not what the traditions say, but how writers from the eighteenth to the twentieth century reread them and produced new revelations about beings from other worlds. Its main result is that almost every ``ancient NHI'' reading in popular circulation can be traced to a dated first appearance in modern print, usually with a documented chain of borrowing from Theosophy through pulp fiction to the contactees and on to the ancient-astronaut authors---the genealogy argued by \textcite{colavito2005cult}, here confirmed at the level of specific wording (Figure~\ref{fig:timeline}; Table~\ref{tab:firsts}). ``First attestation'' means first in the fetched corpus or in the cited scholarship; earlier instances may exist, and proto-forms appear in nineteenth-century speculative writing \parencite{bulwerlytton1871coming,donnelly1882atlantis}.

\begin{figure}[p]
\centering
\includegraphics[width=\textwidth]{figures/fig5_timeline.pdf}
\caption{First appearances in print of ``ancient NHI'' readings. Filled circles: the dated first-attestation table of the modern-reception workstream; open squares: dated first appearances documented by the other regional workstreams. The shaded band marks the period in which the technological reading (ships, rockets, aliens) enters the record.}\label{fig:timeline}
\end{figure}

\begin{table}[p]
\centering\footnotesize
\caption{First attestations of modern ``ancient NHI'' readings and what the source tradition says (after the modern-reception workstream).}\label{tab:firsts}
\begin{tabularx}{\textwidth}{@{}p{4.4cm}p{3.2cm}X@{}}
\toprule
Modern reading & First attestation found & What the source tradition says \\
\midrule
Genesis 6 ``sons of God'' = elementals & Villars 1670 (MOD-05) & Angelic or divine beings (Gen 6:1--4; 1~En 6--16) \\
Divine beings take bodies and interbreed to create humankind & Newbrough 1882 (MOD-08/09) & Enochic Watchers produce giants, not humanity \\
Lords of the Flame incarnate; ``Dzyan'' & Blavatsky 1888 (MOD-10--12) & No premodern witness; admitted paraphrase \\
Culture-bringers from Venus & Scott-Elliot 1904; Leadbeater 1912 & Sanat Kumāra is a Purāṇic mind-born son of Brahmā; no Venus origin \\
Descent ``18 million years ago'' & Bailey 1922 (MOD-22) & No traditional source \\
Venus beings arrive in ships & Lovecraft \& Lumley 1935 (MOD-38) & --- (fiction) \\
Theosophical descent as radiant spacecraft & Leslie 1953 (MOD-25) & --- \\
Atlantean air-ships & Scott-Elliot 1896 (MOD-19) & Clairvoyant claim; Plato has no aircraft \\
\emph{Vimāna} = Atlantean/alien aircraft & Leslie 1953 (MOD-26) & Epic \emph{vimāna} are divine palaces or chariots \\
\emph{Mahābhārata} ``nuclear war'' & BBS pastiche 1993 (MOD-41) & Curse-born iron bolt (MOD-42) \\
\emph{Vaimānika Śāstra} as ancient & Childress 1985, ``1875'' (MOD-43) & Dictated 1918--23 (Josyer) \\
Ezekiel = spacecraft & Williamson 1953 (MOD-29) & Theophany; merkavah mysticism \\
Jesus taken up by a saucer & Downing 1968 (MOD-40) & Resurrection and ascension \\
Earth colonised or ``owned'' & Fort 1919 (MOD-36) & --- \\
Old Ones came ``out of the sky'' & Lovecraft 1926 (MOD-37) & Fiction \\
Anunnaki as astronauts; \emph{shem} = rocket & Sitchin 1976 (MOD-47) & \emph{šumu} = ``name''; gods of heaven and underworld \\
Oannes = alien amphibian & Temple 1976 (MOD-49) & Sea-born sage (Berossus, \emph{apkallu}) \\
Nommo in a spinning ship & Usenet 1995 (MOD-45) & Contested ethnography (van Beek 1991) \\
UFOs = fallen angels or jinn & KeelyNet 1990 (MOD-39); Creighton 1988 (MOD-44) & Reverse reinterpretation \\
Jesus from an ``Evolutionary Level'' & Heaven's Gate c.~1996 (MOD-54) & --- \\
\bottomrule
\end{tabularx}
\end{table}

\paragraph{Enlightenment and esoteric roots.} Swedenborg's \emph{Earths in the Universe} (1758) describes visionary conversations, not visitors: \enquote{it has been granted me by the Lord to speak and converse with spirits and angels who are from other earths, with some for a day, with some for a week, and with some for months}~\pid{[MOD-01]} \parencite{swedenborg1758earths}. He is told \enquote{to be informed that the human race is not from one earth only, but from innumerable earths}~\pid{[MOD-02]} and learns that \enquote{The spirits of Mars are the best of all among the spirits who are from the earths of our solar system}~\pid{[MOD-03]}. His planets stop at Saturn, with no Uranus (discovered 1781)---an early instance of a pattern found throughout the corpus: revelation tracks the astronomy of its day. Villars's satirical \emph{Comte de Gabalis} (1670) first recodes the ``sons of God'' of Genesis 6 as elementals, who \enquote{under the name of the Children of Elohim are distinguished from the Children of Men}~\pid{[MOD-05]} \parencite{villars1670}. Newbrough's channelled \emph{Oahspe} (1882) has gods riding \enquote{ships of fire, coursing the firmament}~\pid{[MOD-07]} and angels who \enquote{clothed they themselves, by force of their wills, with flesh and bones}~\pid{[MOD-08]} beside a pre-human race \parencite{newbrough1882oahspe}.

\paragraph{The Venus lineage.} The clearest textual chain in the corpus runs through Theosophy. Blavatsky's Stanzas of Dzyan (1888)---from a book of which no manuscript or independent witness is known, and which she admits paraphrasing (\enquote{The idea and the spirit of the sentence is here given, as a verbal translation would convey very little to the reader.}~\pid{[MOD-12]})---have the Lords of the Flame descend to \emph{incarnate}: \enquote{"IN THESE SHALL WE DWELL," SAID THE LORDS OF THE FLAME.}~\pid{[MOD-11]} \parencite{blavatsky1888sd}. Venus is Earth's occult elder: \enquote{Thus Earth is the adopted child and younger brother of Venus, but its inhabitants are of their own kind.}~\pid{[MOD-13]} There are no vehicles. Scott-Elliot (1904) makes the Venus origin explicit, \enquote{These came from the scheme of evolution which has Venus as its one physical planet.}~\pid{[MOD-16]} and turns the beings into \enquote{rulers, instructors in religion, and teachers of the arts}~\pid{[MOD-17]} \parencite{scottelliot1904lemuria}. Leadbeater (1912) has Adepts \enquote{from the Venus evolution}~\pid{[MOD-20]} \enquote{transferred to our Earth}~\pid{[MOD-20]} \parencite{leadbeater1912textbook}. Bailey (1922) dates the event---\enquote{in the middle of the Lemurian epoch, approximately eighteen million years ago, occurred a great event}~\pid{[MOD-22]}---without naming Venus \parencite{bailey1922initiation}; Ballard (1934) restores it: \enquote{'Lords of The Flame' from Venus}~\pid{[MOD-24]} \parencite{ballard1934unveiled}. The technological form appears first in fiction, in a Lovecraft revision for William Lumley written in 1935: \enquote{the lords of Venus came through space in their ships to civilise our planet}~\pid{[MOD-38]} \parencite{lovecraft1938typer}. Finally Desmond Leslie's foreword to \emph{Flying Saucers Have Landed} (1953) retells the Theosophical descent as a spacecraft:

\begin{quote}\small About eighteen million years ago, say the strange and ancient legends of our little planet, at a time when Mars, Venus and Earth were in close conjunction, along a magnetic path so formed came a huge, shining, radiant vessel of dazzling power and beauty\par\hfill{\footnotesize --- Flying Saucers Have Landed, Foreword, Foreword, p. [5] of reconstructed PDF \pid{[MOD-25]}}\end{quote}

Its footnote cites only Theosophical authors; Adamski's Venusian ``Orthon'' in the same book gives the doctrine a living informant \parencite{leslieadamski1953landed}. Across this chain reincarnating spiritual hierarchs turn into spacecraft-borne extraterrestrials---in fiction first (1935), in contactee writing afterwards (1953).

\paragraph{Air-ships and vimānas.} Scott-Elliot (1896) claimed that Atlantean air-ships were a realised fact, and compared them to \enquote{the air-ship or flying-machine which Keely in America, and Maxim in this country are now attempting to produce}~\pid{[MOD-19]} \parencite{scottelliot1896atlantis}. Leslie (1953) quotes the passage and inserts the Sanskrit word: \enquote{This is exactly the principle on which the Atlantean vimanas were said to operate.}~\pid{[MOD-27]} This is, as far as the study could find, the first junction of Atlantean clairvoyant air-boats with Indian \emph{vimāna}. BBS files later manufactured a ``nuclear'' Mahābhārata verse, \enquote{hurled a single projectile (rocket) charged with the power of the Universe (nuclear device).}~\pid{[MOD-41]} \parencite{hartman1993ourpast}, while the genuine Mausala Parva has \enquote{an iron bolt through which all the individuals in the race of the Vrishnis and the Andhakas became consumed into ashes.}~\pid{[MOD-42]} \parencite{ganguli1883mahabharata}. Childress (1985) dated the \emph{Vaimānika Śāstra}'s ``rediscovery'' to \enquote{In 1875}~\pid{[MOD-43]} \parencite{childress_aircraft}; the text was dictated in the twentieth century (Section~\ref{sec:asia}).

\paragraph{Fortean writing, pulp and the contactees.} Fort (1919) proposed \enquote{that other worlds explored and colonized here, and fought among themselves for possession}~\pid{[MOD-36]} \parencite{fort1919damned}; Lovecraft's Great Old Ones \enquote{came to the young world out of the sky}~\pid{[MOD-37]} \parencite{lovecraft1928cthulhu}. After Kenneth Arnold's sighting of June 1947, systematic UFO religion emerged, drawing heavily on Theosophy and recasting Ascended Masters as space brothers \parencite{partridge2003ufo}. Williamson (1953) read the Bible accordingly: \enquote{(3) Space visitors, coming to Earth in space craft, are mentioned in the Holy Bible, ancient mythology, and in other documents thousands of years old. These visitors have been coming to Earth for several million years!}~\pid{[MOD-28]} He found in Ezekiel \enquote{a perfect description of Flying Saucers, their landing, and their occupants getting out.}~\pid{[MOD-29]}, and a \enquote{"Great Influx" of the late 1800's}~\pid{[MOD-30]} of hybrid children \parencite{williamson1953tongues}. Ashtar, first channelled in 1952, is \enquote{Commander of ten million Space Men, now occupying bases established within range of your planet}~\pid{[MOD-34]} \parencite{hill1957days,vantassel1952ashtar,helland2003ashtar}; Angelucci's space brother tells of \enquote{another planet in your solar system, the fairest and most radiant of all the planets}~\pid{[MOD-31]} \parencite{angelucci1955secret}. Heaven's Gate completes the process, reading Christian soteriology as \emph{interpretatio technologica}: \enquote{I came to Earth some 2,000 years ago from another physical, biological Evolutionary Level}~\pid{[MOD-54]} \parencite{heavensgate1996jews,richter2012traces}.

\paragraph{The ancient-astronaut synthesis.} Pauwels and Bergier's \emph{Le Matin des magiciens} (1960), which explicitly used Lovecraft, led to Charroux (from 1963) and von Däniken (1968), and then to Sitchin and Temple (both 1976) \parencite{pauwels1960matin,charroux1963histoire,daniken1968erinnerungen,colavito2005cult}. Temple writes that \enquote{a group of alien amphibious beings were credited by the Babylonians with having founded their civilization.}~\pid{[MOD-49]} \parencite{temple1976sirius}; a 1995 Usenet post has Nommo come down \enquote{to earth in a spinning ship.}~\pid{[MOD-45]} \parencite{dolan1995dogon}. Evangelical writers inverted the reading, making UFOs the fallen angels of Genesis 6, and Creighton recast them as jinn \enquote{some of the JINNS could be fully physical and what we call extraterrestrials}~\pid{[MOD-44]} \parencite{keelynet1990nephilim,creighton1988jinns,partridge2004demonology}.

\paragraph{A psychological control.} Flournoy's study of the medium Hélène Smith is a control case of detailed, sincerely reported otherworld revelation. Her ``Martian'' language, he found, \enquote{Martian is only disguised French.}~\pid{[MOD-51]} He explained her romances through \enquote{the phenomena of cryptomnesia, the awakening and setting to work of forgotten memories}~\pid{[MOD-52]}, and noted that her séances took place when \enquote{the conversation more than once turned in the direction of the habitability of Mars, to which the discovery of the famous "canals" has for some years specially directed the attention of the general public}~\pid{[MOD-53]} \parencite{flournoy1900india,henry1901martien}. The contactees' Venusians were plausible before Mariner~2 (1962); afterwards Venusians gave way to Pleiadeans and Zeta Reticulans.

\subsection{Case studies}

The regional sections have already examined the main cases in detail. Here they are summarised side by side (Table~\ref{tab:cases}) to show a recurring structure: a genuine older text or tradition; a modern reframing with a datable first appearance; and a specialist test that the reframing fails.

\begin{table}[ht]
\centering\footnotesize
\caption{Case studies in the construction of ``ancient NHI'' readings.}\label{tab:cases}
\begin{tabularx}{\textwidth}{@{}p{2.3cm}XXX@{}}
\toprule
Case & What the source says & Modern reframing & Specialist test \\
\midrule
Watchers / Anunnaki & Gen 6; 1~En 6--16: angels descend, beget giants, teach arts; Atrahasis: junior gods strike, man made from clay and divine blood & Sitchin 1976: Nefilim as ``People of the Fiery Rockets'', Nibiru astronauts, genetic engineering & Heiser: \emph{šem}/\emph{šumu} = ``name''; \emph{nēberu} a crossing; no Assyriologist endorses Sitchin; apkallu model explains the Watchers (Kilmer, Annus) \\
Dogon Sirius & Griaule \& Dieterlen 1950: \emph{po tolo}, smallest, heaviest star; no contact claim & Temple 1976: Nommo from the Sirius system & van Beek 1991: knowledge unknown outside informant circle; ``twinkling'' Digitaria is visible; contamination possible (Ridpath, Sagan) \\
\emph{Vaimānika Śāstra} & Dictated by Subbaraya Shastry, 1918--23 (or 1903--18) & ``fourth century B.C.'', ``rediscovered in 1875'' (Childress 1985) & Mukunda et al. 1974: modern Sanskrit, craft unworkable, drawings by a local draughtsman \\
\emph{Kebra Nagast} & Wagons raised a cubit by Michael; Solomon's air-vessel; ``chariot'' as palladium & Flying machines, anti-gravity & Solomon-and-the-wind legend cycle; chs. 113--117 no flight \\
Bep Kororoti & Métraux 1960; Posey 1983: storm-shaman ascends in lightning & von Däniken 1973: spacesuited teacher & Peret's gloss and editorial ``(or the future? - Red.)'' inside indigenous speech \\
Hopi Ant People & Voth 1905; Lockett 1933: emergence through lower planes; kachinas as ancestral spirits & Waters 1963; Blumrich 1979: Ant People, star kachinas & Geertz 1983/1994: White Bear uninitiated; syncretic authorship \\
Wandjina & Ancestral beings who painted their own images & von Däniken: helmeted astronauts & Frederick \& O'Connor 2009: reading descends from colonial denial of authorship \\
Chitauri & Imbulu: tailed shape-shifter of a false-bride tale (Theal 1886) & Mutwa 1999 (via Icke): reptilian Chitauri, Greys as servants & Chidester 2002; Podolecka 2023: ``not Zulu myths'' \\
Magonia & Agobard c.~815: popular belief in cloud-ships of weather-magicians & Villars 1670: abduction by sylphs; Vallée 1969: fairy lore = UFO lore & Chain traced; Vallée's quotation of Evans-Wentz differs from the 1911 text (to be checked) \\
\bottomrule
\end{tabularx}
\end{table}

\paragraph{Two further observations.} First, the Ezekiel case shows the process at work within a single career: Blumrich began by trying to refute von Däniken---\enquote{with the intention of proving von Daniken wrong.}~\pid{[MOD-50]}---and ended by building an engineering model of the prophet's ``spaceship'' \parencite{blumrich1974ezekiel}. Second, feedback loops are common. Mutwa's Nommo comes from Temple; the Hopi ``Anunnaki'' from Sitchin; the \emph{Vaimānika}'s Atlantean resonance from Leslie. Fringe readings return, sometimes through Indigenous or non-Western intermediaries, as ``tradition''---which is why dating the first appearance of each element matters.

%=====================================================================
\section{Discussion}\label{sec:discussion}
%=====================================================================

\subsection{Evaluating the hypotheses motif by motif}

Table~\ref{tab:hyp} summarises the evaluation. The central finding can be stated simply: \emph{no passage in the corpus requires H1}. That is not a proof that contact never happened; it is the finding that the attested motifs are fully accounted for by other, independently evidenced processes. The strongest candidates for a ``visitor'' reading---Callaway's white being who descends in a fog, Chinigchinich ``from the stars'', the Kayapó people who descend from the sky, the \emph{Kebra Nagast}'s aerial wagons, Ezekiel's wheels---each have simpler readings within their own traditions, and in several cases (Callaway, the \emph{Kebra Nagast}, Ezekiel) the text itself supplies them.

H1 also fails its most discriminating prediction. If the texts recorded visitors from advanced civilisations, one would expect at least occasional information beyond the authors' horizon. What one finds instead is the astronomy and technology of each text's own time: the 364-day schematic year and gated flat earth of 1~Enoch; Swedenborg's planets ending at Saturn; Hélène Smith's Martian built on French; contactee Venusians before Mariner~2; Scott-Elliot's Atlantean air-ships modelled on Keely and Maxim. The one celebrated counter-example, Dogon knowledge of Sirius~B, dissolves under ethnographic restudy.

\begin{table}[ht]
\centering\footnotesize
\caption{Evaluation of the five hypotheses by motif (P/E evidence in the five traditional areas). ``Strong'' = the hypothesis explains the attested features with independent support; ``part'' = explains some cases; ``--'' = little or no role. H1 is never \emph{required}.}\label{tab:hyp}
\begin{tabularx}{\textwidth}{@{}lXccccc@{}}
\toprule
Motif & Where strong in P/E evidence & H1 & H2 & H3 & H4 & H5 \\
\midrule
SKY & All areas; dynastic charters (Japan, Korea, Ife) & not req. & strong & strong & part & part \\
CULT & Near East, Americas, Celtic, Yoruba & not req. & strong & part & part & part \\
FORB & Near East only (Popol Vuh inverts) & not req. & strong & -- & -- & part \\
HYB & Near East, Celtic, Asia & not req. & strong & strong & part & strong \\
VEH & Asia (epic), Near East, Kebra Nagast & not req. & strong & part & strong & strong \\
STAR & Americas; /Xam; 1 Enoch & not req. & part & strong & part & strong \\
TIME & Celtic; Viṣṇu Purāṇa & not req. & part & strong & -- & part \\
ABD & Celtic; all areas & not req. & part & strong & part & part \\
PLUR & Asia (metaphysical) & not req. & part & strong & -- & strong \\
DIM & Celtic, Americas, Near East, Asia & not req. & part & strong & part & part \\
CHAN & (modern only) & not req. & -- & part & -- & strong \\
WAR & Near East, Asia & not req. & strong & part & -- & part \\
\bottomrule
\end{tabularx}
\end{table}

\paragraph{H2, diffusion and literary borrowing,} is decisively supported where names, plots and specific details are shared along known routes: the Watchers complex (Mesopotamia → Second Temple Judaism → Ethiopia, Manichaeism and Islam); the Ethiopic Watchers and Solomon's wind-borne travel in the \emph{Kebra Nagast}; the Northeast Asian heavenly-descent kingship complex; Chinese Penglai in Urashima; the Yavana association of flying \emph{yantras}; the Star Husband and Earth-diver complexes mapped by Thompson; and, in the modern period, the chain from Theosophy to the contactees.

\paragraph{H3, shared cognitive and structural myth-making,} best explains structural similarity without shared detail: layered cosmoi reached by rope, tree, spider-thread or rainbow bridge; sky-marriage and false-bride tale-types; hidden peoples attached to landscape features on both sides of the Atlantic; the global ``group of youths/sisters'' Pleiades pattern; and otherworld time lapse. Cognitive accounts---intuitive attribution of agency to sky phenomena, minimally counterintuitive spirits---explain why these motifs spread easily, though not the specific shared names and plots that H2 explains.

\paragraph{H4, translation and editorial artefacts,} is demonstrable repeatedly: Cory's ``destitute of reason'' against the variant ``endowed''; Hewitt's ``man-beings''; Hirata's gloss in the \emph{Kojiki}; Griffith's ``animated vessels''; Josyer's ``Aeroplanes'' and ``Aeronautics''; Platt's ``Rapid Transit''; Gregory's ``through the air and the high air''; Kingsley's missionary crops; Callaway's leading questions; Wyndham's blank verse; Griaule's co-construction; the Pleyades editor's ``(or the future? - Red.)''.

\paragraph{H5, modern reinterpretation,} accounts for almost every explicitly extraterrestrial element in the popular picture, with datable first appearances between 1670 and 1999 and a concentration of the technological reading between 1896 and 1953 (Figure~\ref{fig:timeline}).

\subsection{What remains genuinely interesting}

An evidence-first study should say what is left unexplained, not only what is debunked. Four problems remain genuinely interesting.

\paragraph{The universality of the otherworld taking.} Abduction into an otherworld, with return, is attested in every traditional area: the Thunder's bride, Enoch's summons, Raivata's visit to Brahmā, Urashima, Bran and Oisín, the fairy-taken nurse, the Star Husband's wives. Diffusion can explain particular clusters, but not the whole distribution. Why so many cultures imagined that the invisible neighbours \emph{take people} deserves an explanation that neither contact nor mere borrowing supplies.

\paragraph{The lapse of time.} The supernatural lapse of time is attested in Ireland by the seventh or eighth century, in the Purāṇa in the first millennium CE, in Walter Map in the twelfth century, and in Japan in the eighth-century \emph{fudoki} tradition. Independent invention through a shared schema of the otherworld as the land of the dead or the ever-living is the most economical hypothesis, but, as the Celtic workstream noted, no phylogenetic tale-type analysis has been done for this motif.

\paragraph{Bullard's uniformity argument.} Bullard examined about 300 abduction reports and found a stable sequence---capture, examination, return---in 163 of 193 codable cases \parencite{bullard1987,bullard1989}. He read abductions as the supernatural kidnap narrative in technological guise, supporting the structural parallel with fairy lore; but he also stressed that abduction narratives are \emph{more uniform} than folklore under ordinary oral dynamics should be, and called their stability a genuine puzzle for the folklorist. Critics have suggested that hypnotic regression and investigator filtering produce the uniformity. The question is empirical and could be tested by comparing the variability of fairy legends (e.g. migratory legend types ML~5070 and ML~4075) with that of abduction reports.

\paragraph{Experience-centred explanations.} Hufford's study of the ``Old Hag'' (sleep paralysis) shows how some night-visitor and abduction features can be cross-culturally stable without either contact or diffusion \parencite{hufford1982}; Bullard conceded its relevance. The Celtic comparison table of the workstream suggests that sourceless light is the single most specific shared feature between fairy halls and abduction settings, while vehicles, protection rituals and social functions differ. Scholars such as \textcite{kripal2010authors} and \textcite{pasulka2019cosmic} argue that the \emph{experiences} reported deserve study as religious phenomena whatever their cause; nothing in this study contradicts that. What the study does show is that the \emph{texts} usually cited as ancient evidence are not the right evidence for such experiences.

\subsection{Limitations}

\begin{itemize}
\item \textbf{Colonial and missionary mediation.} Most ethnographic texts were collected by outsiders under colonial conditions, from few informants, and translated through outsider categories. They are documented versions, not transparent indigenous voice.
\item \textbf{Outdated translations.} Much of the corpus rests on superseded Victorian and Edwardian translations (Charles, Budge, Griffith, Chamberlain, King, Kramer 1944, Rogers). Key passages should be checked against modern critical editions before being cited as evidence of what a tradition said; some philological questions (e.g. the Ge`ez behind Budge's ``cubit'' and ``air (or winds)'') remain open.
\item \textbf{English-language bias.} The corpus is overwhelmingly English; a few French, Portuguese and Spanish texts were read. Sources in the original languages were consulted only through translators' notes.
\item \textbf{Sampling of the site.} Sacred-texts is a curated library with a strong Anglo-Celtic and Theosophical skew and thin coverage of Africa, Islam beyond the Qur'ān and Sufism, China, and living oral traditions. The motif counts measure the shape of a sampled corpus, not the frequency of motifs in the traditions.
\item \textbf{No fieldwork.} The study is entirely textual; it cannot speak to contemporary belief in any community.
\item \textbf{Aggregator reliability.} Bibliotecapleyades pages include unauthorised reproductions, unattributed translations (the Genesis Apocryphon, the Tan'gun text), mislabelling (Korean texts filed under ``China''), altered numbers (in Rutt's Tongmyŏng), and editorial interpolations. Quotations from Sitchin, Temple, Blumrich, Leslie and Vallée taken from such copies are flagged for checking against print.
\item \textbf{Unverified bibliographic details.} Some bibliography entries were marked by researchers as uncertain (e.g. the attribution of the sacred-texts \emph{Vaimānika} preface to J.~B. Hare, the probable identification of the Tan'gun translation with Peter H. Lee, the volume of Creighton's Dropa investigation). These are kept with explanatory notes; no missing detail has been supplied.
\end{itemize}

%=====================================================================
\section{Conclusion}\label{sec:conclusion}
%=====================================================================

The world's sacred and folk literatures, as represented in the English-accessible corpus examined here, are full of non-human intelligences: angels, Watchers and jinn; orisha and heaven-folk; devas, nāgas and gandharvas; kami; the \emph{síd} peoples; star-persons, kachinas and Wandjina. They descend from heaven, teach, marry, beget, take people away, and fight. These motifs are genuine, ancient in several cases, and in some cases---sky descent, the otherworld journey, hidden peoples---nearly universal.

What the texts do not contain is equally clear. No traditional text names a planet of origin. No primary vehicle is a machine built by non-humans; vehicles are divine gifts, magical properties, visions or human automata. Forbidden knowledge from heaven is a single Near Eastern literary tradition. Channelled contact is modern. The technological reading---ships, rockets, astronauts, hybrids produced by procedure---enters the record between 1896 and 1953, through Theosophy, pulp fiction and the contactees, and it reweights the traditional motifs rather than merely reading them.

The evidence therefore supports a layered explanation. Diffusion and literary borrowing (H2) explain the shared names and plots; shared cognitive and structural myth-making (H3) explains the widespread structures; translation and editorial mediation (H4) explain many of the most ``modern-looking'' details; and modern reinterpretation (H5) explains almost every explicitly extraterrestrial element. Literal contact (H1) cannot be excluded by textual evidence, and the paper makes no claim to have excluded it; but none of the motifs, as attested, requires it. What remains genuinely open is not whether ancient texts describe spacecraft---they do not---but why human cultures so consistently imagine invisible neighbours who take people away, and why time runs differently where they live.

\section*{Acknowledgements and data statement}
This paper synthesises the work of six regional research workstreams and two indexing passes conducted for Douglas Oliveira in September 2026. All passages, codes and ratings are available in the files described in Appendix~\ref{app:data}.

\printbibheading[heading=bibintoc,title={References}]
\begin{multicols}{2}\raggedright
\printbibliography[heading=none]
\end{multicols}

\appendix
%=====================================================================
\section{Motif codebook}\label{app:codebook}
%=====================================================================
\begin{longtable}{@{}lp{9.2cm}r@{}}
\caption{Motif codebook (definitions from \texttt{index.md}) and number of tags in the corpus.}\label{tab:codebook}\\
\toprule Code & Definition & Tags \\ \midrule \endfirsthead
\toprule Code & Definition & Tags \\ \midrule \endhead
\bottomrule \endfoot
\textsc{sky} & Sky descent: gods/beings descend from sky or stars to earth & 92 \\
\textsc{cult} & Culture-bringer: non-human teacher gives arts, law, agriculture, writing, kingship & 67 \\
\textsc{forb} & Forbidden knowledge transfer (Watchers/Azazel, Harut \& Marut, Prometheus-type theft) & 20 \\
\textsc{hyb} & Hybrid offspring / divine-human unions / embryo transfer / changelings & 55 \\
\textsc{veh} & Celestial vehicles: flying chariots, wagons, vimanas, sky-boats, 'airships', saucers & 74 \\
\textsc{star} & Star-origin ancestry \& star-people (Pleiades, Sirius, Morning Star, Venus) & 37 \\
\textsc{time} & Otherworld time dilation (years pass in hours; return to dust) & 10 \\
\textsc{abd} & Abduction / otherworld journey / ascent to sky-world and return & 60 \\
\textsc{plur} & Plurality of inhabited worlds / beings from other planets & 36 \\
\textsc{dim} & Other planes/dimensions: fairy, elemental, subterranean or parallel hidden peoples & 65 \\
\textsc{chan} & Channeled / contactee communication with non-human intelligences & 28 \\
\textsc{war} & War or rebellion among sky beings; fallen/expelled beings & 34 \\
\end{longtable}
\noindent\emph{Coding notes.} Codes are etic and descriptive; the emic identity (angel, jinn, orisha, deva, \emph{síd}-folk, kachina, etc.) is recorded in each passages notes. A passage may carry several codes. Source types: P primary text in translation; E ethnographic record; S secondary scholarship or compilation; F fringe, esoteric or new-revelation literature; C critical scholarship on a fringe claim. Strength: strong = motif unambiguous in the passage; moderate = clear but partial; weak = requires interpretation or is a negative/control case.

%=====================================================================
\section{Key passages}\label{app:passages}
%=====================================================================
The table lists every corpus passage quoted in this paper, with tradition, text, motif codes, source type (P/E/S/F/C) and strength (S strong, M moderate, W weak). Full quotations, URLs, locators and interpretive notes are in \texttt{evidence/corpus.json}.

\scriptsize
\setlength{\tabcolsep}{2.5pt}
\renewcommand{\arraystretch}{0.95}
\begin{longtable}{@{}p{1.6cm}p{3.6cm}p{6.9cm}p{2.1cm}cc@{}}
\caption{Passages quoted in this paper (210 of 314 in the corpus).}\label{tab:passages}\\
\toprule ID & Tradition & Text & Motifs & Type & Str. \\ \midrule \endfirsthead
\toprule ID & Tradition & Text & Motifs & Type & Str. \\ \midrule \endhead
\bottomrule \endfoot
\multicolumn{6}{@{}l}{\textbf{Near East \& Abrahamic}} \\[2pt]
\texttt{NE01} & Second Temple Judaism (Enochic) & 1 Enoch (Book of Watchers) & SKY HYB & P & S \\
\texttt{NE03} & Second Temple Judaism (Enochic) & 1 Enoch (Book of Watchers) & FORB CULT & P & S \\
\texttt{NE04} & Second Temple Judaism (Enochic) & 1 Enoch (Book of Watchers) & FORB & P & S \\
\texttt{NE11} & Second Temple Judaism (Enochic) & 1 Enoch (Book of Watchers) & STAR WAR & P & S \\
\texttt{NE13} & Second Temple Judaism (Enochic) & 1 Enoch (Similitudes) & FORB CULT & P & M \\
\texttt{NE16} & Second Temple Judaism (Enochic) & 1 Enoch (Book of Noah fragment) & HYB & P & S \\
\texttt{NE17} & Second Temple Judaism (Jubilees) & Book of Jubilees & SKY CULT ABD & P & S \\
\texttt{NE22} & Jewish/Christian pseudepigrapha (2 Enoch) & 2 Enoch (Book of the Secrets of Enoch) & VEH & P & W \\
\texttt{NE25} & Jewish/Christian pseudepigrapha (2 Enoch) & 2 Enoch (Book of the Secrets of Enoch) & ABD & P & M \\
\texttt{NE26} & Hebrew Bible & Genesis 6 & HYB SKY & P & S \\
\texttt{NE27} & Hebrew Bible & Genesis 5 & ABD & P & M \\
\texttt{NE28} & Hebrew Bible & Ezekiel 1 & VEH SKY & P & S \\
\texttt{NE29} & Hebrew Bible & Ezekiel 10 & VEH & P & S \\
\texttt{NE30} & Manichaeism (Enochic reception, Central… & Book of Giants (Manichaean fragments, Kawan) & SKY FORB WAR & P & S \\
\texttt{NE33} & Second Temple Judaism (Qumran) & Book of Giants (Qumran: 4Q531, 4Q530) & WAR HYB & P & S \\
\texttt{NE39} & Ancient Mesopotamia & Berossus, Babyloniaca (via Alexander Polyhistor, in… & CULT DIM & P & S \\
\texttt{NE41} & Ancient Mesopotamia & Sumerian Mythology (Inanna and Enki) & CULT & S & S \\
\texttt{NE43} & Ancient Mesopotamia & Atrahasis, Tablet I (OB version) & WAR & P & S \\
\texttt{NE44} & Ancient Mesopotamia & Atrahasis, Tablet I (OB version) & HYB CULT & P & S \\
\texttt{NE46} & Ancient Mesopotamia & Adapa (Amarna fragment and NA copies) & ABD FORB CULT & P & S \\
\texttt{NE49} & Islam & Qur'an, Sura 2 (al-Baqara) & FORB SKY & P & S \\
\texttt{NE50} & Islam & Qur'an, Sura 72 (al-Jinn) & DIM FORB CHAN & P & S \\
\texttt{NE51} & Islam & Qur'an, Sura 55 (al-Rahman) and Sura 15 & DIM PLUR & P & M \\
\texttt{NE53} & Islam & The Bible, the Koran, and the Talmud (Biblical Legends of the… & DIM & S & M \\
\texttt{NE54} & Modern reinterpretation / critique & The 12th Planet, ch. 5 'The Nefilim: People of the Fiery Rockets' & VEH SKY HYB & F & W \\
\texttt{NE55} & Modern reinterpretation / critique & Sitchin's Mesopotamian Rocket Ships: How to Ignore Ancient… & VEH & C & S \\
\texttt{NE57} & Ancient Mesopotamia & Enuma Elish, Tablet VI & CULT HYB & P & M \\
\multicolumn{6}{@{}l}{\textbf{Africa}} \\[2pt]
\texttt{AFR-KN-01} & Ethiopian Orthodox (Solomonic royal epic) & Kebra Nagast (The Queen of Sheba and her only son Menyelek) & VEH & P & M \\
\texttt{AFR-KN-02} & Ethiopian Orthodox (Solomonic royal epic) & Kebra Nagast (The Queen of Sheba and her only son Menyelek) & VEH SKY & P & S \\
\texttt{AFR-KN-03} & Ethiopian Orthodox (Solomonic royal epic) & Kebra Nagast (The Queen of Sheba and her only son Menyelek) & VEH & P & S \\
\texttt{AFR-KN-05} & Ethiopian Orthodox (Solomonic royal epic) & Kebra Nagast (The Queen of Sheba and her only son Menyelek) & SKY WAR & P & S \\
\texttt{AFR-KN-06} & Ethiopian Orthodox (Solomonic royal epic) & Kebra Nagast (The Queen of Sheba and her only son Menyelek) & HYB & P & S \\
\texttt{AFR-KN-09} & Ethiopian Orthodox (Solomonic royal epic) & Kebra Nagast (The Queen of Sheba and her only son Menyelek) & VEH & P & W \\
\texttt{AFR-MLB-03} & Bantu-speaking peoples (various) & Myths and Legends of the Bantu & SKY & S & S \\
\texttt{AFR-MLB-04} & Bantu-speaking peoples (various) & Myths and Legends of the Bantu & ABD HYB & S & S \\
\texttt{AFR-MLB-06} & Bantu-speaking peoples (various) & Myths and Legends of the Bantu & ABD & S & S \\
\texttt{AFR-MLB-07} & Bantu-speaking peoples (various) & Myths and Legends of the Bantu & ABD DIM & S & S \\
\texttt{AFR-MLB-08} & Bantu-speaking peoples (various) & Myths and Legends of the Bantu & SKY HYB & S & S \\
\texttt{AFR-MLB-09} & Bantu-speaking peoples (various) & Myths and Legends of the Bantu & SKY CULT & S & S \\
\texttt{AFR-MLB-11} & Bantu-speaking peoples (various) & Myths and Legends of the Bantu & ABD & S & S \\
\texttt{AFR-MLB-12} & Bantu-speaking peoples (various) & Myths and Legends of the Bantu & ABD CULT PLUR & S & S \\
\texttt{AFR-RSA-01} & Zulu (amaZulu) & The Religious System of the Amazulu & SKY & E & S \\
\texttt{AFR-RSA-02} & Zulu (amaZulu) & The Religious System of the Amazulu & SKY DIM & E & M \\
\texttt{AFR-RSA-03} & Zulu (amaZulu) & The Religious System of the Amazulu & SKY & E & M \\
\texttt{AFR-IFE-01} & Yoruba (Ife) & Myths of Ife & SKY CULT & E & S \\
\texttt{AFR-IFE-02} & Yoruba (Ife) & Myths of Ife & CULT & E & S \\
\texttt{AFR-IFE-03} & Yoruba (Ife) & Myths of Ife & WAR & E & M \\
\texttt{AFR-YOR-01} & Yoruba & The Yoruba-Speaking Peoples of the Slave Coast of West Africa & ABD SKY & E & M \\
\texttt{AFR-YL-01} & Yoruba & Yoruba Legends & CULT & E & W \\
\texttt{AFR-FJO-01} & Kongo (Fjort/Vili, Loango coast) & Notes on the Folklore of the Fjort & ABD CULT & E & S \\
\texttt{AFR-FJO-02} & Kongo (Fjort/Vili, Loango coast) & Notes on the Folklore of the Fjort & CULT & S & M \\
\texttt{AFR-SBF-01} & /Xam San (Bushman), Northern Cape & Specimens of Bushman Folklore & STAR & E & M \\
\texttt{AFR-SBF-03} & /Xam San (Bushman), Northern Cape & Specimens of Bushman Folklore & STAR & E & S \\
\texttt{AFR-XFT-01} & Xhosa & Kaffir Folk-Lore & DIM & E & M \\
\texttt{AFR-DAS-01} & Gullah/Geechee (African American,… & Drums and Shadows & ABD & E & M \\
\texttt{AFR-DOG-01} & Dogon (Mali) & "Un système soudanais de Sirius" (Griaule \& Dieterlen 1950),… & STAR & E & W \\
\texttt{AFR-DOG-02} & Dogon (Mali) & "Un système soudanais de Sirius" (Griaule \& Dieterlen 1950),… & STAR & E & W \\
\texttt{AFR-DOG-03} & Dogon (Mali) & "Un système soudanais de Sirius" (Griaule \& Dieterlen 1950),… & STAR & E & W \\
\texttt{AFR-DOG-05} & Dogon (Mali) & "Un système soudanais de Sirius" (Griaule \& Dieterlen 1950),… & CULT & E & W \\
\texttt{AFR-DOG-06} & Dogon (modern reinterpretation) & The Sirius Mystery & STAR CULT & F & W \\
\texttt{AFR-DOG-09} & Dogon (critique) & Dogon Restudied: A Field Evaluation of the Work of Marcel Griaule & STAR & C & W \\
\texttt{AFR-DOG-10} & Dogon (critique) & Dogon Restudied: A Field Evaluation of the Work of Marcel Griaule & STAR & C & W \\
\texttt{AFR-DOG-12} & Dogon (critique) & "Investigating the Sirius 'Mystery'", Skeptical Inquirer 3(1) & STAR & C & W \\
\texttt{AFR-MUT-01} & Zulu / "African" (modern informant) & Interview "Great Zulu Shaman and Elder Credo Mutwa - On Alien… & STAR PLUR & F & W \\
\texttt{AFR-MUT-02} & Zulu / "African" (modern informant) & Interview "Great Zulu Shaman and Elder Credo Mutwa - On Alien… & WAR HYB SKY & F & W \\
\texttt{AFR-MUT-03} & Zulu / "African" (modern informant) & Interview "Great Zulu Shaman and Elder Credo Mutwa - On Alien… & ABD CHAN & F & W \\
\texttt{AFR-MUT-04} & Zulu / "African" (modern informant) & Interview "Great Zulu Shaman and Elder Credo Mutwa - On Alien… & DIM HYB & F & W \\
\texttt{AFR-MUT-07} & Zulu / "African" (modern informant) & Credo Mutwa, first-person essays ("The Origins of the Gods",… & VEH CULT & F & W \\
\texttt{AFR-POD-01} & Zulu (scholarly assessment of Mutwa) & "History of the Bantu Peoples by Credo Mutwa - Zulu Tradition or… & CHAN & C & W \\
\multicolumn{6}{@{}l}{\textbf{South \& East Asia}} \\[2pt]
\texttt{asia-001} & Hindu (Sanskrit epic) & Mahabharata, Vana Parva (Book 3) & VEH WAR & P & S \\
\texttt{asia-002} & Hindu (Sanskrit epic) & Mahabharata, Vana Parva (Book 3) & VEH WAR & P & S \\
\texttt{asia-004} & Hindu (Sanskrit epic) & Mahabharata, Vana Parva (Book 3), Indralokagamana Parva & VEH SKY ABD & P & S \\
\texttt{asia-005} & Hindu (Sanskrit epic) & Mahabharata, Vana Parva (Book 3), Indralokagamana Parva & STAR ABD PLUR & P & M \\
\texttt{asia-006} & Hindu (Sanskrit epic) & Mahabharata, Vana Parva (Book 3), Nivatakavacha-yuddha Parva & VEH WAR DIM & P & S \\
\texttt{asia-007} & Hindu (Sanskrit epic) & Mahabharata, Vana Parva (Book 3), Nivatakavacha-yuddha Parva & VEH WAR & P & S \\
\texttt{asia-009} & Hindu (Sanskrit epic) & Mahabharata, Adi Parva (Book 1), Sambhava Parva & HYB SKY & P & S \\
\texttt{asia-011} & Hindu (Sanskrit epic) & Mahabharata, Mausala Parva (Book 16) & WAR & P & S \\
\texttt{asia-012} & Hindu (Sanskrit epic) & Ramayan of Valmiki, Book V (Sundara Kanda) & VEH & P & S \\
\texttt{asia-013} & Hindu (Sanskrit epic) & Ramayan of Valmiki, Book VI (Yuddha Kanda) & VEH & P & S \\
\texttt{asia-015} & Vedic (Rigveda) & Rig Veda 4.36 (Hymn to the Rbhus) & VEH CULT & P & M \\
\texttt{asia-016} & Vedic (Rigveda) & Rig Veda 1.116 (Hymn to the Asvins) & VEH ABD & P & M \\
\texttt{asia-017} & Hindu (Purana) & Vishnu Purana, Book IV, ch. 1 & TIME ABD & P & S \\
\texttt{asia-018} & Hindu (Purana) & Vishnu Purana, Book IV, ch. 1 & TIME ABD & P & S \\
\texttt{asia-019} & Hindu (Purana) - translator commentary & Vishnu Purana, Book IV, ch. 1, Wilson note 33 & TIME & S & M \\
\texttt{asia-022} & Hindu (Purana) & Vishnu Purana, Book II, ch. 7 & PLUR & P & S \\
\texttt{asia-023} & Jain (Svetambara canon) & Kalpa Sutra: Life of Mahavira, Lecture 2 (SBE 22) & SKY VEH & P & M \\
\texttt{asia-025} & Buddhist (Pali canon) & Brahmajala Sutta (Digha Nikaya 1), Dialogues of the Buddha I & DIM PLUR SKY & P & M \\
\texttt{asia-027} & Buddhist (Mahayana) & Saddharma-pundarika (Lotus Sutra), ch. 1 (SBE 21) & PLUR & P & S \\
\texttt{asia-028} & Buddhist (Mahayana) & Saddharma-pundarika (Lotus Sutra), ch. 14 (SBE 21) & PLUR SKY & P & M \\
\texttt{asia-029} & Modern Hindu (20th-c. channelled… & Vymaanika-Shaastra (Vaimanika Shastra), English title page & CHAN VEH & P & S \\
\texttt{asia-030} & Modern Hindu (20th-c. channelled… & Vymaanika-Shaastra, First Chapter & VEH CHAN PLUR & P & S \\
\texttt{asia-031} & Modern Hindu - institutional study… & Vymanika Shastra Rediscovered, ch. 1 Background & CHAN & S & S \\
\texttt{asia-032} & Scientific critique (IISc) & A critical study of the work "Vymanika Shastra" (Scientific Opinion… & CHAN & C & S \\
\texttt{asia-033} & Scientific critique (IISc) & A critical study of the work "Vymanika Shastra" (Scientific Opinion… & CHAN VEH & C & S \\
\texttt{asia-034} & Scientific critique (IISc) & A critical study of the work "Vymanika Shastra" (Scientific Opinion… & VEH & C & M \\
\texttt{asia-035} & Editorial commentary (online edition) & The Vymaanika-Shaastra: Preface to the sacred-texts edition & CHAN & C & M \\
\texttt{asia-036} & Hindu (medieval shilpa-shastra) via… & Samarangana Sutradhara of Bhoja, ch. 31 (Yantravidhana), vv. 95-98,… & VEH & S & S \\
\texttt{asia-037} & Hindu/Buddhist literature via scholarly… & Brhatkathaslokasamgraha (Budhasvamin) episode, as summarised by V.… & VEH & S & M \\
\texttt{asia-040} & Shinto (Japanese court myth-history) & Kojiki, Section XXXIV: The August Reign in Himuka & SKY CULT & P & S \\
\texttt{asia-041} & Shinto - translator commentary & Kojiki, Section XXXIV, Chamberlain note 4 & SKY VEH & S & S \\
\texttt{asia-043} & Shinto (Japanese court myth-history) & Nihongi, Book III (Emperor Jimmu), Part 2 & SKY VEH HYB & P & S \\
\texttt{asia-046} & Korean foundation myth & Tan'gun legend (Samguk yusa 1, Kogi/"Old Record"), as reproduced on… & SKY CULT HYB & P & S \\
\texttt{asia-047} & Korean foundation myth & Lay of King Tongmyong (Tongmyong-wang p'yon), as reproduced on… & VEH SKY HYB & P & S \\
\texttt{asia-049} & Tibetan Buddhism - scholarly critique & Mistaken Foreign Myths about Shambhala & DIM ABD & C & M \\
\texttt{asia-050} & Modern ancient-astronaut literature… & Shambhala and the UFO Connection (compilation) quoting D.H.… & WAR VEH & F & W \\
\texttt{asia-051} & Modern ancient-astronaut legend… & Los Dropa (Fate Magazine, Aug. 1998) & VEH HYB STAR & F & W \\
\multicolumn{6}{@{}l}{\textbf{Celtic \& N. Europe}} \\[2pt]
\texttt{celtic-01} & Irish (medieval pseudo-history) & Lebor Gabála Érenn (Book of Invasions), Tuatha Dé Danann section & SKY & P & M \\
\texttt{celtic-02} & Irish (medieval pseudo-history) & Lebor Gabála Érenn, verse on the coming of the Tuatha Dé Danann & SKY DIM & P & S \\
\texttt{celtic-03} & Irish (medieval, via early-20th-c.… & The Fairy-Faith in Celtic Countries, ch. IV 'The People of the… & SKY CULT & S & M \\
\texttt{celtic-05} & Irish (medieval saga) & Cath Maige Tuired (Second Battle of Mag Tuired) & SKY & P & S \\
\texttt{celtic-06} & Irish (medieval saga) & Cath Maige Tuired & CULT FORB & P & M \\
\texttt{celtic-09} & Irish (medieval saga) & Cath Maige Tuired & CULT & P & S \\
\texttt{celtic-12} & Irish (Revival retelling of medieval… & Gods and Fighting Men, Part I Bk I 'Fight with the Firbolgs' & SKY & P & M \\
\texttt{celtic-14} & Irish (Revival retelling) & Gods and Fighting Men, Part I Bk IV 'Bodb Dearg' & DIM & P & S \\
\texttt{celtic-15} & Irish (Revival retelling of Echtrae… & Gods and Fighting Men, Part I Bk IV 'Call to Connla' & ABD DIM VEH & P & S \\
\texttt{celtic-17} & Irish (Old Irish immram/echtrae) & Immram Brain (Voyage of Bran) & DIM ABD & P & S \\
\texttt{celtic-19} & Irish (Old Irish immram/echtrae) & Immram Brain (Voyage of Bran) & TIME & P & S \\
\texttt{celtic-21} & Irish (Fenian cycle; early-modern lay… & Gods and Fighting Men, 'Oisin's Story' & TIME & P & S \\
\texttt{celtic-22} & Scottish Highland (early modern) & The Secret Common-Wealth, ch. 1 'Of the Subterranean Inhabitants' & DIM & E & S \\
\texttt{celtic-23} & Scottish Highland (early modern) & The Secret Common-Wealth, ch. 2 & DIM & E & M \\
\texttt{celtic-24} & Scottish Highland (early modern) & The Secret Common-Wealth, ch. 4 & ABD HYB & E & S \\
\texttt{celtic-26} & Scottish Highland (early modern) & The Secret Common-Wealth, ch. 15 & ABD DIM & E & S \\
\texttt{celtic-28} & Scottish (Border ballad) & Thomas Rymer (Child 37A) & ABD DIM & P & S \\
\texttt{celtic-29} & Scottish (Border ballad) & Tam Lin (Child 39A) & ABD HYB & P & S \\
\texttt{celtic-31} & Irish (modern field testimony, Co. Sligo) & The Fairy-Faith in Celtic Countries, ch. II 'Ireland': testimony of… & DIM ABD & E & S \\
\texttt{celtic-32} & Irish (modern field testimony, Co. Sligo) & The Fairy-Faith in Celtic Countries, ch. II: testimony of Patrick… & PLUR SKY DIM & E & W \\
\texttt{celtic-34} & Irish (Revival visionary; AE) & The Fairy-Faith in Celtic Countries, 'An Irish Mystic's Testimony' & DIM CHAN & E & W \\
\texttt{celtic-36} & Scottish Gaelic (modern field… & The Fairy-Faith in Celtic Countries, ch. II 'Scotland' & TIME ABD & E & S \\
\texttt{celtic-37} & Scottish Gaelic (modern field… & The Fairy-Faith in Celtic Countries, testimony of Murdoch MacLean,… & HYB & E & M \\
\texttt{celtic-38} & Welsh (modern field testimony) & The Fairy-Faith in Celtic Countries, ch. II 'Wales' & TIME ABD & E & S \\
\texttt{celtic-40} & Celtic (early-20th-c. scholarly… & The Fairy-Faith in Celtic Countries, ch. XI 'Science and Fairies' & DIM & S & M \\
\texttt{celtic-41} & Anglo-Welsh / European (medieval… & The Science of Fairy Tales, ch. VII 'The Supernatural Lapse of Time… & TIME ABD & S & S \\
\texttt{celtic-42} & British (19th-c. comparative folklore) & The Science of Fairy Tales, ch. V 'Changelings' & HYB & S & S \\
\texttt{celtic-44} & Norse (Snorri's mythography) & Gylfaginning (Prose Edda), ch. XVII & DIM & P & M \\
\texttt{celtic-46} & Carolingian (primary historical source) & Agobard of Lyon, De grandine et tonitruis (On Hail and Thunder), ch. 2 & VEH & P & S \\
\texttt{celtic-47} & French esoteric (Rosicrucian… & Comte de Gabalis, Discourse V & VEH DIM & F & M \\
\texttt{celtic-48} & French esoteric & Comte de Gabalis, Discourse V & ABD VEH & F & M \\
\texttt{celtic-51} & Modern ufology (comparative thesis) & Passport to Magonia, chapter on sexual contact (Pleyades excerpt) & HYB & F & M \\
\texttt{celtic-52} & Modern ufology (argument paper) & Five Arguments Against the Extraterrestrial Origin of Unidentified… & ABD DIM VEH & F & W \\
\multicolumn{6}{@{}l}{\textbf{Americas \& Pacific}} \\[2pt]
\texttt{AP01} & Timagami Ojibwa (Anishinaabe),… & Stith Thompson (ed.), Tales of the North American Indians (No. L,… & STAR ABD HYB & E & S \\
\texttt{AP03} & Pan-North American (comparative) & Stith Thompson (ed.), Tales of the North American Indians,… & STAR ABD & S & S \\
\texttt{AP04} & Seneca (Haudenosaunee/Iroquois) & Stith Thompson (ed.), Tales of the North American Indians (No. V,… & SKY DIM & E & S \\
\texttt{AP05} & Onondaga (Haudenosaunee) & J.N.B. Hewitt, Iroquoian Cosmology, First Part (BAE 21st Annual… & SKY DIM & E & S \\
\texttt{AP06} & Onondaga (Haudenosaunee) & J.N.B. Hewitt, Iroquoian Cosmology, First Part (BAE 21st Annual… & SKY & E & S \\
\texttt{AP07} & Haudenosaunee (Iroquoian), editorial & J.N.B. Hewitt, Iroquoian Cosmology, First Part (BAE 21st Annual… & DIM & S & M \\
\texttt{AP08} & Cherokee (Tsalagi), Southeast & James Mooney, Myths of the Cherokee (BAE 19th Annual Report, Pt. 1)… & STAR & E & S \\
\texttt{AP09} & Cherokee (Tsalagi), Southeast & James Mooney, Myths of the Cherokee (BAE 19th Annual Report, Pt. 1)… & STAR ABD & E & S \\
\texttt{AP10} & Cherokee (Tsalagi), Southeast & James Mooney, Myths of the Cherokee (BAE 19th Annual Report, Pt. 1)… & DIM & E & S \\
\texttt{AP11} & Cherokee (Tsalagi), Southeast & James Mooney, Myths of the Cherokee (BAE 19th Annual Report, Pt. 1)… & DIM & E & S \\
\texttt{AP12} & Cherokee (Tsalagi), Southeast & James Mooney, Myths of the Cherokee (BAE 19th Annual Report, Pt. 1)… & CHAN ABD & E & S \\
\texttt{AP13} & Hopi (Oraibi), Southwest & H.R. Voth, The Traditions of the Hopi (Field Columbian Museum Publ.… & DIM PLUR & E & S \\
\texttt{AP14} & Hopi, Southwest & Hattie Greene Lockett, The Unwritten Literature of the Hopi (Univ.… & DIM & E & S \\
\texttt{AP15} & Hopi, Southwest & Hattie Greene Lockett, The Unwritten Literature of the Hopi & PLUR DIM & E & S \\
\texttt{AP16} & Hopi (modern reception; fringe) & Ellen Lloyd, "Watchers of the Hopi - Astonishing Secrets of… & STAR VEH CULT & F & W \\
\texttt{AP17} & Hopi (modern reception; fringe) & Gary A. David, "The Anthills of Orion - Ancient Star Beings of the… & STAR DIM & F & W \\
\texttt{AP21} & Acjachemen (Juaneño), coastal southern… & Gerónimo Boscana, Chinigchinich (tr. Alfred Robinson in Life in… & CULT STAR SKY & E & S \\
\texttt{AP22} & Acjachemen (Juaneño; Serrano narrative) & Gerónimo Boscana, Chinigchinich (tr. Alfred Robinson in Life in… & SKY CULT & E & S \\
\texttt{AP23} & K'iche' (Quiché) Maya, highland Guatemala & Popol Vuh: The Sacred Book of the Ancient Quiché Maya (Preamble) & CULT & P & M \\
\texttt{AP24} & K'iche' Maya & Popol Vuh: The Sacred Book of the Ancient Quiché Maya (Part III,… & FORB & P & S \\
\texttt{AP25} & K'iche' Maya & Popol Vuh: The Sacred Book of the Ancient Quiché Maya (Part III,… & FORB & P & S \\
\texttt{AP26} & K'iche' Maya & Popol Vuh: The Sacred Book of the Ancient Quiché Maya (Part I, Ch. 7) & STAR & P & S \\
\texttt{AP29} & Yucatec Maya (reported by a Spanish… & Diego de Landa, Relación de las cosas de Yucatán (as Yucatan Before… & CULT & P & S \\
\texttt{AP31} & Nahua (modern reception; fringe) & Arjun Walia, "Was Quetzalcoatl an Extraterrestrial?" (Collective… & CULT SKY & F & W \\
\texttt{AP32} & Inca / Andean (Collao, Cuzco) & Pedro Sarmiento de Gamboa, History of the Incas (Historia Índica),… & CULT & P & M \\
\texttt{AP34} & Kayapó (Mebêngôkre), Kuben-kran-kegn… & Alfred Métraux, "Mythes et contes des Indiens Cayapo (groupe… & SKY STAR & E & S \\
\texttt{AP35} & Kayapó (Mebêngôkre), Kuben-kran-kegn & Alfred Métraux, "Mythes et contes des Indiens Cayapo", Revista do… & SKY & E & S \\
\texttt{AP36} & Kayapó (Mebêngôkre), Gorotire & Darrell A. Posey, "Keeping of Stingless Bees by the Kayapó Indians… & SKY & E & S \\
\texttt{AP38} & Kayapó (Peret's Portuguese narrative) & João Américo Peret, "Bep-Kororoti", Cadernos Brasileiros de… & SKY VEH & E & W \\
\texttt{AP40} & Kayapó (modern reception; fringe) & Erich von Däniken, In Search of Ancient Gods (1973/1974), as… & SKY & F & W \\
\texttt{AP44} & Māori (per Smith / Whatahoro) & S. Percy Smith, The Lore of the Whare-wananga, Part I: Te… & PLUR & E & W \\
\texttt{AP46} & Euahlayi (Yuwaalaraay) & K. Langloh Parker, The Euahlayi Tribe & STAR & E & S \\
\texttt{AP47} & Kimberley (Worrorra, Ngarinyin,… & Ursula Frederick \& Sue O'Connor, "Wandjina, graffiti and heritage",… & CULT DIM & C & S \\
\texttt{AP48} & Kimberley Wandjina (history of reception) & Frederick \& O'Connor 2009, section "The Wandjina mystique" (quoting… & SKY VEH & C & S \\
\texttt{AP49} & Kimberley Wandjina (modern reception;… & Anon., "Los Wandjinas y El Tiempo de los Sueños" (Spanish web page) & SKY CULT & F & W \\
\multicolumn{6}{@{}l}{\textbf{Modern reception}} \\[2pt]
\texttt{MOD-01} & Swedenborgian (18th-c. visionary… & Earths in the Universe (De Telluribus in Mundo Nostro Solari) & PLUR CHAN & F & S \\
\texttt{MOD-02} & Swedenborgian & Earths in the Universe & PLUR & F & S \\
\texttt{MOD-03} & Swedenborgian & Earths in the Universe & PLUR CHAN & F & S \\
\texttt{MOD-05} & Early-modern Rosicrucian/Paracelsian… & Comte de Gabalis, Discourse II & HYB DIM FORB & F & S \\
\texttt{MOD-07} & Faithism (Oahspe), American Spiritualism & Oahspe, Book of Aph, Son of Jehovih, ch. III & VEH SKY WAR & F & S \\
\texttt{MOD-08} & Faithism (Oahspe) & Oahspe, Book of Jehovih, ch. VI & HYB SKY FORB & F & S \\
\texttt{MOD-11} & Theosophy (Blavatsky) & The Secret Doctrine, vol. II, Stanza VII.24 (cont.) & SKY CULT & F & S \\
\texttt{MOD-12} & Theosophy (Blavatsky) & The Secret Doctrine, vol. II, Stanza VII.24 footnote & CHAN & F & S \\
\texttt{MOD-13} & Theosophy (Blavatsky) & The Secret Doctrine, vol. II, "Parent Stars and Sister Planets" & PLUR STAR & F & S \\
\texttt{MOD-16} & Theosophy (second generation) & The Lost Lemuria, "Teachers of the Lemurian Race" & SKY CULT PLUR & F & S \\
\texttt{MOD-17} & Theosophy (second generation) & The Lost Lemuria, "Teachers of the Lemurian Race" & CULT SKY & F & S \\
\texttt{MOD-19} & Theosophy (second generation) & The Story of Atlantis, "Air-Ships" & VEH & F & S \\
\texttt{MOD-20} & Theosophy (Adyar, Leadbeater) & A Textbook of Theosophy, ch. IX "The Planetary Chains" & SKY CULT PLUR & F & S \\
\texttt{MOD-22} & Theosophy (Bailey / Arcane School) & Initiation, Human and Solar, ch. IV "The Founding of the Hierarchy" & SKY CULT CHAN & F & S \\
\texttt{MOD-24} & I AM Religious Activity (Ballard) & Unveiled Mysteries, ch. IX "Venus Visits the Royal Teton" & SKY CULT PLUR CHAN & F & S \\
\texttt{MOD-25} & 1950s contactee literature… & Flying Saucers Have Landed, Foreword (Book One by Leslie) & SKY VEH CULT PLUR & F & S \\
\texttt{MOD-27} & 1950s contactee literature (Leslie) & Flying Saucers Have Landed, ch. 7 & VEH & F & S \\
\texttt{MOD-28} & 1950s contactee literature (Williamson) & Other Tongues--Other Flesh, Prologue & SKY VEH PLUR & F & S \\
\texttt{MOD-29} & 1950s contactee literature (Williamson) & Other Tongues--Other Flesh, ch. "The Prophets" & VEH SKY & F & S \\
\texttt{MOD-30} & 1950s contactee literature (Williamson) & Other Tongues--Other Flesh, ch. "The Wanderers" & HYB STAR CHAN & F & S \\
\texttt{MOD-31} & 1950s contactee literature (Angelucci) & The Secret of the Saucers, ch. III "My Meeting With Neptune" & WAR STAR PLUR & F & S \\
\texttt{MOD-34} & Ashtar Command (channelled) & In Days To Come, Introduction & CHAN PLUR VEH & F & S \\
\texttt{MOD-36} & Fortean anomalism & The Book of the Damned, ch. 12 & PLUR WAR VEH & F & S \\
\texttt{MOD-37} & Weird fiction (Lovecraft) as source of… & "The Call of Cthulhu" & SKY STAR CHAN & P & S \\
\texttt{MOD-38} & Weird fiction (Lovecraft \& Lumley) & "The Diary of Alonzo Typer" & SKY VEH CULT PLUR & P & S \\
\texttt{MOD-41} & 1990s BBS ufology (vimana/nuclear-war) & "Flying Aircraft and Nuclear War... of the Past" OURPAST.ASC & VEH WAR & F & W \\
\texttt{MOD-42} & Hindu epic (control text) & The Mahabharata, Mausala Parva, section 1 & WAR & P & S \\
\texttt{MOD-43} & Ancient-astronaut literature (vimana) & "Ancient Indian Aircraft Technology" AIAC & VEH & F & W \\
\texttt{MOD-44} & 1990s BBS ufology (Islamic reframing) & "The True Nature of UFO Entities" JINNS.ASC & DIM PLUR & F & M \\
\texttt{MOD-45} & 1990s Usenet ufology (Dogon/Sirius) & "The Dogon" DOGON.UFO & STAR VEH CULT SKY & F & W \\
\texttt{MOD-49} & Ancient-astronaut literature (Temple) & The Sirius Mystery, ch. 9 "A Fable" & CULT STAR SKY & F & M \\
\texttt{MOD-50} & Ancient-astronaut literature (Blumrich) & The Spaceships of Ezekiel (excerpt/intro) & VEH SKY & F & M \\
\texttt{MOD-51} & Psychology of mediumship (Flournoy case… & From India to the Planet Mars, ch. VI "The Martian Language" & CHAN PLUR & S & S \\
\texttt{MOD-52} & Psychology of mediumship (Flournoy case… & From India to the Planet Mars, ch. XI "Conclusion" & CHAN & S & S \\
\texttt{MOD-53} & Psychology of mediumship (Flournoy case… & From India to the Planet Mars, ch. V "The Martian Cycle" & PLUR CHAN & S & S \\
\texttt{MOD-54} & Heaven's Gate (UFO new religious movement) & "The Jews and Christians Promote Lies - Unknowingly" & SKY VEH CHAN WAR & F & S \\
\end{longtable}\normalsize

%=====================================================================
\section{Data and code availability}\label{app:data}
%=====================================================================
All files are in the project directory \texttt{nhi-research/} and in the deliverable package.\begingroup\raggedright
\begin{description}[leftmargin=1.5em,style=nextline]
\item[\texttt{index.md}, \texttt{index.json}] Topic index of sacred-texts.com: 77 categories, 261 entries, motif codebook, gaps and caveats (corrected; original kept as \texttt{index.md.bak\_before\_corrections}).
\item[\texttt{index\_pleyades.md}, \texttt{index\_pleyades.json}] Supplementary index of 166 Biblioteca Pleyades items with the gap table G1--G9.
\item[\texttt{evidence/<area>.json}, \texttt{\_notes.md}, \texttt{.bib}] For each of the six areas: verified passages, analytic synthesis and bibliography.
\item[\texttt{evidence/corpus.json}] The merged corpus of 314 passages with an added \texttt{area} field and, where relevant, an \texttt{editor\_flag} recording a known correction or a check still required.
\item[\texttt{paper/references.bib}] Merged, deduplicated bibliography (10 duplicate works under different keys merged; key map in \texttt{paper/data/bib\_aliases.json}).
\item[\texttt{paper/data/}] Motif~$\times$~area matrices (all passages; P/E; strong P/E; fringe) and source-type/strength breakdowns as CSV and JSON.
\item[\texttt{paper/scripts/}] \texttt{fix\_index.py} (index corrections), \texttt{build\_data.py} (merge, matrices, bibliography), \texttt{make\_figures.py} (figures), \texttt{build\_tex.py} (inserts verbatim quotations from \texttt{corpus.json} into \texttt{main\_src.tex} to produce \texttt{main.tex}, and checks citation keys).
\item[\texttt{paper/figures/}] Figures 1--6 as PDF and PNG.
\end{description}\endgroup

\begin{figure}[ht]
\centering
\includegraphics[width=0.8\textwidth]{figures/fig6_strength.pdf}
\caption{Strength ratings by area.}\label{fig:strength}
\end{figure}

\end{document}
